\documentclass[12pt]{article}
\usepackage[spanish]{babel}
\usepackage[T1]{fontenc}
\usepackage[utf8]{inputenc}
\usepackage{lmodern}
\usepackage{amsmath,amssymb}
\usepackage{amsthm}
\usepackage{geometry}
\geometry{margin=2.5cm}
\usepackage{comment}
\renewcommand{\familydefault}{\sfdefault}
\usepackage{mathrsfs}
\usepackage{graphicx}
\usepackage{float}
\usepackage{subcaption}
\usepackage{booktabs,tabularx,array}
\usepackage{enumitem}
\usepackage{titlesec}
\usepackage[authoryear]{natbib}
\bibliographystyle{plainnat}
\usepackage[hidelinks]{hyperref}

\newcommand{\fuente}[1]{\textnormal{(#1)}}
\newcommand{\fuentecite}[2]{\textnormal{(#1; \citealp{#2})}}


\newtheorem{definicion}{Definición}[section]
\newtheorem{lema}{Lema}[section]
\newtheorem{proposicion}{Proposición}[section]
\newtheorem{teorema}{Teorema}[section]
\newtheorem{corolario}{Corolario}[section]
\newtheorem{observacion}{Observación}[section]
\newtheorem{ejemplo}{Ejemplo}[section]
\newtheorem{conjetura}{Conjetura}[section]
\newtheorem{problema}{Problema abierto}[section]

\title{Control Informacional\\
\large Proyecto de Investigación II --- ProInt\_II}

\author{
\textbf{Alejandro Cruz López} \\
Universidad Autónoma Metropolitana \\
Unidad Iztapalapa
}

\date{
\vspace{1em}
\textbf{Asesores} \\[0.3em]
Dr.\ Asael Fabián Martínez Martínez \\
Dr.\ Joaquín Delgado Fernández
\vspace{1em}
26 de mayo de 2026
}


\begin{document}
\maketitle

\section*{Introducción general}

El estudio de distribuciones de probabilidad finitas suele abordarse desde dos perspectivas principales: como objetos estáticos sujetos a estimación, o como estados que se actualizan mediante reglas inferenciales o dinámicas específicas. En distintos contextos, sin embargo, aparecen secuencias de distribuciones cuya evolución constituye por sí misma el fenómeno de interés. Esto conduce a una pregunta distinta: ¿Cómo describir transformaciones entre estados probabilísticos preservando su estructura intrínseca, sin reducir necesariamente el análisis a un modelo estocástico externo?.

En el caso finito, toda distribución de probabilidad puede representarse como un punto del símplex. Esta representación aparece de manera natural en la teoría axiomática de la probabilidad (\cite{Kolmogorov1933}). A partir de esta observación, el problema del cambio informacional puede reformularse como el estudio de transformaciones internas del símplex compatibles con la interpretación probabilística. En este trabajo se fija un dominio preciso ---estados informacionales discretos inducidos por particiones medibles finitas--- y se introducen transformaciones elementales definidas mediante correcciones relativas. Estas transformaciones dan lugar a operadores bien definidos sobre el símplex y permiten incorporar funcionales informacionales para comparar estados (\cite{Csiszar1975}).

El punto de vista adoptado no pretende reemplazar marcos ya establecidos para dinámicas estocásticas o reglas particulares de actualización. Un ejemplo clásico de tales marcos lo proporcionan las cadenas de Markov. También existen formulaciones variacionales para la evolución de medidas en espacios métricos (\cite{JordanKinderlehrerOtto1998}). Más modestamente, el objetivo aquí es aislar un formalismo matemático en el que las transformaciones entre distribuciones finitas puedan estudiarse directamente como objeto de análisis. En esta formulación, los cambios entre estados se modelan mediante correcciones relativas que preservan la masa total. En el desarrollo posterior, estas correcciones admitirán interpretación composicional (\cite{Aitchison1986}). También admitirán una interpretación por cambio de medida en el caso discreto. Cuando sea pertinente, se indicará además su interpretación geométrica local en términos de la métrica de Fisher (\cite{Amari2016}).

El interés de este problema no es únicamente formal. En diversas situaciones aparecen familias finitas de probabilidades o vectores normalizados cuya evolución contiene información relevante sobre estabilidad, convergencia o respuesta al entorno. En particular, las actualizaciones multiplicativas sobre el símplex aparecen en optimización y aprendizaje (\cite{KivinenWarmuth1997}). Por otra parte, la estructura geométrica del símplex ha sido estudiada sistemáticamente en análisis composicional (\cite{Aitchison1986}). También ha sido analizada desde la geometría de la información (\cite{AmariNagaoka2000}).

El objetivo central de este manuscrito unificado es construir y organizar un marco básico para el estudio de transformaciones deterministas de estados probabilísticos finitos. En particular, se busca:
\begin{enumerate}
    \item fijar con precisión el dominio probabilístico discreto en el que se trabajará;
    \item introducir correcciones relativas y operadores asociados que preserven la estructura del símplex;
    \item relacionar dichas transformaciones con cambios discretos de medida;
    \item incorporar funcionales informacionales y estudiar la dinámica que inducen en el caso discreto.
\end{enumerate}

El alcance del trabajo es deliberadamente acotado. Se considera exclusivamente el caso de distribuciones finitas, formuladas como puntos del símplex de probabilidad, y se estudian las estructuras mínimas necesarias para pasar de estados aislados a transformaciones y dinámicas entre ellos. En este sentido, el documento tiene un carácter fundacional: no pretende agotar el problema general, sino establecer con rigor el núcleo matemático que servirá de base para desarrollos posteriores.

La organización del texto es la siguiente. En la Sección A se fija el dominio probabilístico y su representación geométrica en el símplex. En la Sección B se introducen correcciones locales que preservan la estructura probabilística y se muestra su interpretación como cambios discretos de medida. En la Sección C se construyen operadores globales asociados a dichas correcciones y se estudia su estructura algebraica. Finalmente, en la Sección D se incorporan funcionales informacionales, en particular la divergencia de Kullback--Leibler (\cite{KullbackLeibler1951}). Sobre esta base se analiza una dinámica determinista sobre el interior del símplex con interpretación geométrica de tipo Fisher (\cite{Amari2016}).

\subsection*{Relación con marcos de referencia}

El desarrollo del trabajo utiliza herramientas provenientes de varios marcos que actúan sobre el mismo dominio, el símplex de probabilidades $\Delta_n$, pero que cumplen funciones distintas dentro de la exposición.

\begin{itemize}
    \item \textbf{Análisis composicional (Aitchison).} Proporciona la estructura algebraico--geométrica interna del símplex cuando las transformaciones relevantes son multiplicativas y seguidas de cierre o normalización. En particular, las operaciones de perturbación y potenciación permiten trabajar con cantidades relativas dentro del símplex (\cite{Aitchison1986}).

    \item \textbf{Geometría de divergencias (Csiszár).} Aporta funcionales de comparación entre distribuciones. En particular, las $f$-divergencias proporcionan un marco general para medir discrepancias entre estados probabilísticos (\cite{Csiszar1975}). Dentro de esta familia, la divergencia de Kullback--Leibler ocupa un lugar central en este trabajo (\cite{KullbackLeibler1951}).

    \item \textbf{Geometría de la información (Amari).} Proporciona la métrica de Fisher y el lenguaje diferencial que permite interpretar gradientes y direcciones de cambio en el interior del símplex (\cite{AmariNagaoka2000}). Este punto de vista se utilizará sólo cuando sea necesario para interpretar localmente la dinámica (\cite{Amari2016}).

    \item \textbf{Formulaciones variacionales en espacios de medidas (Wasserstein/JKO).} Este marco no se desarrolla en el documento como teoría principal. Sin embargo, se toma como referencia cuando se discuten analogías entre funcionales, esquemas discretos y dinámicas de descenso (\cite{JordanKinderlehrerOtto1998}). La formulación abstracta de estos flujos en espacios métricos puede consultarse en (\cite{AmbrosioGigliSavare2005}).
\end{itemize}

La convención adoptada será la siguiente. Las transformaciones elementales se formulan principalmente en la geometría composicional del símplex (\cite{Aitchison1986}). Los funcionales de comparación se expresan mediante divergencias informacionales (\cite{Csiszar1975}). Cuando resulte necesario, las interpretaciones diferenciales se harán explícitas en términos de la métrica de Fisher (\cite{Amari2016}). De manera análoga, las referencias a esquemas variacionales se entenderán sólo como punto de comparación conceptual (\cite{JordanKinderlehrerOtto1998}).
\section{Sección A — Probabilidad y símplex}

Este sección  fija el dominio probabilístico discreto en el que se desarrollará el trabajo. El punto de partida es un espacio de probabilidad en el sentido de Kolmogórov (\cite{Kolmogorov1933}). A partir de una partición medible finita se asocia a la medida un vector de probabilidades, que se interpretará como estado informacional discreto. Dicho vector pertenece naturalmente al símplex de probabilidad.

\begin{observacion}
Conviene distinguir desde el inicio entre la medida $P$ sobre $(\Omega,\mathcal F)$ y el estado discreto $X$ inducido por una partición $\Pi$. Un mismo espacio de probabilidad puede dar lugar a distintos estados al variar la partición considerada.
\end{observacion}

\begin{definicion}[Espacio de probabilidad]\label{def:espacio_prob}
Un \emph{espacio de probabilidad} es una terna $(\Omega,\mathcal F,P)$ tal que:
\begin{enumerate}
    \item $\Omega$ es un conjunto no vacío;
    \item $\mathcal F$ es una $\sigma$-álgebra de subconjuntos de $\Omega$;
    \item $P:\mathcal F\to[0,1]$ es una medida con $P(\Omega)=1$.
\end{enumerate}
\end{definicion}

\begin{definicion}[Partición medible finita]\label{def:particion}
Sea $(\Omega,\mathcal F,P)$ un espacio de probabilidad. Una \emph{partición medible finita} de $\Omega$ es una familia finita
\[
\Pi=\{E_1,\dots,E_n\}\subset \mathcal F,\qquad n\ge 1,
\]
tal que:
\begin{enumerate}
    \item $E_i\cap E_j=\varnothing$ si $i\neq j$;
    \item $\bigcup_{i=1}^n E_i=\Omega$.
\end{enumerate}

Diremos que $\Pi$ es \emph{no degenerada} si además
\[
P(E_i)>0,\qquad i=1,\dots,n.
\]
\end{definicion}

\begin{observacion}
La hipótesis de no degeneración no forma parte de la definición general de partición, pero será útil cuando intervengan cocientes o logaritmos.
\end{observacion}

\begin{definicion}[Símplex de probabilidad e interior]\label{def:simplex}
Para $n\ge 1$ se define el \emph{símplex de probabilidad} por
\[
\Delta_n:=\Bigl\{x\in\mathbb R^n:\ x_i\ge 0,\ \sum_{i=1}^n x_i=1\Bigr\}.
\]
Su interior relativo en el hiperplano $\sum_{i=1}^n x_i=1$ está dado por
\[
\Delta_n^\circ:=\{x\in\Delta_n:\ x_i>0\ \text{para todo } i\}.
\]
\end{definicion}

\begin{observacion}
El conjunto $\Delta_n$ será el espacio general de estados. Cuando se requiera positividad estricta, por ejemplo en expresiones con cocientes o logaritmos, se trabajará en $\Delta_n^\circ$.
\end{observacion}

\begin{definicion}[Estado informacional discreto]\label{def:estado}
Sea $(\Omega,\mathcal F,P)$ un espacio de probabilidad y sea $\Pi=\{E_1,\dots,E_n\}$ una partición medible finita. El \emph{estado informacional discreto} asociado a $\Pi$ es el vector
\[
X=(X_1,\dots,X_n)\in\mathbb R^n,
\qquad
X_i:=P(E_i),\quad i=1,\dots,n.
\]
\end{definicion}

\begin{proposicion}\label{prop:estado_simplex}
Sea $X$ el estado informacional discreto asociado a una partición medible finita $\Pi=\{E_1,\dots,E_n\}$.
Entonces:
\begin{enumerate}
    \item $X\in\Delta_n$;
    \item si $\Pi$ es no degenerada, entonces $X\in\Delta_n^\circ$.
\end{enumerate}
\end{proposicion}

\begin{proof}
Como $E_i\in\mathcal F$, se tiene $X_i=P(E_i)\ge 0$ para todo $i$. Además, por disjunción de la partición y aditividad numerable de $P$,
\[
\sum_{i=1}^n X_i
=
\sum_{i=1}^n P(E_i)
=
P\Bigl(\bigcup_{i=1}^n E_i\Bigr)
=
P(\Omega)
=
1.
\]
Por tanto, $X\in\Delta_n$. Si además $P(E_i)>0$ para todo $i$, entonces $X_i>0$ para todo $i$, y en consecuencia $X\in\Delta_n^\circ$.
\end{proof}

\begin{observacion}
El paso de $(\Omega,\mathcal F,P)$ a $X$ depende de la partición elegida. En consecuencia, $X$ no representa a la medida $P$ en toda su estructura, sino sólo su distribución sobre las clases de observación determinadas por $\Pi$.
\end{observacion}

\begin{definicion}[Operador de cierre]\label{def:cierre}
Sea
\[
\mathbb R_{>0}^n:=\{z=(z_1,\dots,z_n)\in\mathbb R^n:\ z_i>0\ \text{para todo } i\}.
\]
El \emph{operador de cierre} es la aplicación
\[
\mathscr{C}:\mathbb R_{>0}^n\to\Delta_n^\circ,
\qquad
\mathscr{C}(z):=\frac{1}{\sum_{i=1}^n z_i}\,z.
\]
\end{definicion}

\begin{observacion}
Para todo $\lambda>0$ y todo $z\in\mathbb R_{>0}^n$ se cumple
\[
\mathscr{C}(\lambda z)=\mathscr{C}(z).
\]
Esta invarianza de escala expresa que el cierre conserva únicamente la información relativa entre coordenadas y además la normalización coincide con la operación de cierre utilizada en análisis composicional para pasar de vectores positivos a composiciones sobre el símplex (\cite{Aitchison1986}).
\end{observacion}

\begin{observacion}
Cuando se requiera interpretar transformaciones multiplicativas internas del símplex, se utilizará el formalismo composicional introducido por Aitchison (\cite{Aitchison1986}). En particular, el interior del símplex admite una estructura algebraica natural mediante perturbación y potenciación. Para los resultados de esa teoría que sean necesarios más adelante, se remitirá a la bibliografía correspondiente.
\end{observacion}
\section{Sección B — Correcciones locales y cambio de medida}

En este sección  se introduce el mecanismo elemental de transformación sobre el símplex. Dado un estado $X\in\Delta_n$, se definen correcciones relativas admisibles que preservan la no negatividad y la normalización, se asocia a cada una de ellas un operador local sobre el espacio de estados y se muestra que dicha actualización coincide con un cambio de medida discreto sobre la partición subyacente.

\begin{definicion}[Conjunto de correcciones admisibles en $X$]\label{def:CX}
Sea $X=(X_1,\dots,X_n)\in\Delta_n$. Se define el \emph{conjunto de correcciones admisibles en $X$} por
\[
C(X)
:=
\Bigl\{
Y=(Y_1,\dots,Y_n)\in\mathbb{R}^n
\;\Big|\;
X_i(1+Y_i)\ge 0\ \text{para todo } i,
\ \sum_{i=1}^n X_i(1+Y_i)=1
\Bigr\}.
\]
\end{definicion}

\begin{observacion}\label{obs:neutralidad}
Si $Y\in C(X)$, entonces
\[
\sum_{i=1}^n X_iY_i=0.
\]
En efecto, como $\sum_{i=1}^n X_i=1$ y $\sum_{i=1}^n X_i(1+Y_i)=1$, se obtiene
\[
\sum_{i=1}^n X_iY_i
=
\sum_{i=1}^n X_i(1+Y_i)-\sum_{i=1}^n X_i
=
1-1
=
0.
\]
Esta identidad expresa la neutralidad de masa de la corrección, es decir, $\sum_{i=1}^n X_iY_i=0$.
\end{observacion}

\begin{definicion}[Operador informacional local]\label{def:Tloc}
Sea $X\in\Delta_n$ y sea $Y\in C(X)$. El \emph{operador informacional local} asociado a $Y$ en $X$ es el vector
\[
T_Y^{\mathrm{loc}}(X)
=
\bigl(T_Y^{\mathrm{loc}}(X)_1,\dots,T_Y^{\mathrm{loc}}(X)_n\bigr)\in\mathbb{R}^n,
\]
definido por
\[
T_Y^{\mathrm{loc}}(X)_i:=X_i(1+Y_i),\qquad i=1,\dots,n.
\]
\end{definicion}

\begin{observacion}\label{obs:frontera}
No es necesario exigir $X\in\Delta_n^\circ$ para definir $T_Y^{\mathrm{loc}}$. Si $X_i=0$, entonces
\[
T_Y^{\mathrm{loc}}(X)_i=0
\]
independientemente del valor de $Y_i$. Por tanto, el operador local está bien definido en todo $\Delta_n$. La restricción al interior $\Delta_n^\circ$ sólo será necesaria más adelante, cuando intervengan expresiones que requieran positividad estricta.
\end{observacion}

\begin{proposicion}[Preservación local del símplex]\label{prop:preservacionLocal}
Sea $X\in\Delta_n$ y sea $Y\in C(X)$. Entonces
\[
T_Y^{\mathrm{loc}}(X)\in\Delta_n.
\]
\end{proposicion}

\begin{proof}
Por definición de $C(X)$,
\[
X_i(1+Y_i)\ge 0,\qquad i=1,\dots,n,
\]
y
\[
\sum_{i=1}^n X_i(1+Y_i)=1.
\]
Como
\[
T_Y^{\mathrm{loc}}(X)_i=X_i(1+Y_i),
\]
se sigue que todas las coordenadas de $T_Y^{\mathrm{loc}}(X)$ son no negativas y que su suma es $1$. Por tanto,
\[
T_Y^{\mathrm{loc}}(X)\in\Delta_n.
\]
\end{proof}

\begin{proposicion}[Cambio de medida discreto inducido por una corrección local]\label{prop:QYprobabilidad}
Sea $(\Omega,\mathcal F,P)$ un espacio de probabilidad y sea
\[
\Pi=\{E_1,\dots,E_n\}\subset\mathcal F
\]
una partición medible finita no degenerada. Defínase
\[
X_i:=P(E_i),\qquad i=1,\dots,n,
\]
de modo que $X=(X_1,\dots,X_n)\in\Delta_n^\circ$.

Sea $Y\in C(X)$ y defínase
\[
h:\Omega\to[0,\infty),
\qquad
h(\omega):=1+Y_i
\quad\text{si }\omega\in E_i.
\]
Sea además
\[
Q_Y:\mathcal F\to[0,\infty],
\qquad
Q_Y(A):=\int_A h(\omega)\,dP(\omega),
\qquad A\in\mathcal F.
\]
Entonces $Q_Y$ es una medida de probabilidad sobre $(\Omega,\mathcal F)$.
\end{proposicion}

\begin{proof}
Como $Y\in C(X)$, se tiene
\[
X_i(1+Y_i)\ge 0,\qquad i=1,\dots,n.
\]
Puesto que $X_i=P(E_i)>0$, se sigue que
\[
1+Y_i\ge 0,\qquad i=1,\dots,n.
\]
Por tanto,
\[
h(\omega)\ge 0,\qquad \omega\in\Omega.
\]
En consecuencia, la aplicación
\[
A\mapsto \int_A h\,dP
\]
define una medida no negativa sobre $(\Omega,\mathcal F)$.

Además, usando que $\Pi=\{E_1,\dots,E_n\}$ es una partición de $\Omega$ y que $h$ es constante en cada sección ,
\[
Q_Y(\Omega)
=
\int_\Omega h\,dP
=
\sum_{i=1}^n \int_{E_i} h\,dP
=
\sum_{i=1}^n (1+Y_i)P(E_i)
=
\sum_{i=1}^n X_i(1+Y_i)
=
1.
\]
Por tanto, $Q_Y$ es una medida de probabilidad.
\end{proof}

\begin{proposicion}[Compatibilidad entre $T_Y^{\mathrm{loc}}$ y el cambio de medida]\label{prop:compatibilidadCambioMedida}
Sea $(\Omega,\mathcal F,P)$ un espacio de probabilidad y sea
\[
\Pi=\{E_1,\dots,E_n\}\subset\mathcal F
\]
una partición medible finita no degenerada. Defínase
\[
X_i:=P(E_i),\qquad i=1,\dots,n,
\]
de modo que $X=(X_1,\dots,X_n)\in\Delta_n^\circ$.

Sea $\omega \in \Omega$. Sea $Y \in C(X)$ y defínase $h:\Omega \to (0,\infty)$ por
\[
h(\omega):=1+Y_i \quad \text{si } \omega \in E_i.
\]
Sea además
\[
Q_Y(A):=\int_A h(\omega)\,dP(\omega),
\qquad A\in\mathcal F.
\]
Entonces, para cada $i=1,\dots,n$, se cumple
\[
Q_Y(E_i)=T_Y^{\mathrm{loc}}(X)_i.
\]
\end{proposicion}

\begin{proof}
Como $h$ es constante sobre cada sección  $E_i$ y vale $1+Y_i$ en dicho sección ,
\[
Q_Y(E_i)
=
\int_{E_i} h(\omega)\,dP(\omega)
=
(1+Y_i)P(E_i)
=
(1+Y_i)X_i.
\]
Por la definición del operador local,
\[
(1+Y_i)X_i=T_Y^{\mathrm{loc}}(X)_i.
\]
Por tanto,
\[
Q_Y(E_i)=T_Y^{\mathrm{loc}}(X)_i.
\]
\end{proof}

\begin{observacion}
La proposición anterior identifica la actualización local con la redistribución de probabilidad inducida por la medida $Q_Y$ sobre la partición original. En este sentido, el operador $T_Y^{\mathrm{loc}}$ no actúa sólo como transformación algebraica sobre el símplex, sino como representación discreta de un cambio de medida por sección s.
\end{observacion}
\section{Sección C — Operadores globales multiplicativos y composición}

En esta sección  se pasa del mecanismo local del sección 2 a una familia de operadores definida de manera uniforme sobre el interior del símplex. La idea es considerar correcciones multiplicativas positivas y renormalizar el resultado para conservar la suma unitaria. Esto da lugar a operadores globales sobre $\Delta_n^\circ$ y permite describir su composición de forma explícita.

\begin{definicion}[Correcciones globales]\label{def:Cglob}
Se define el conjunto de \emph{correcciones globales} por
\[
C_{\mathrm{glob}}
:=
\Bigl\{
Y=(Y_1,\dots,Y_n)\in\mathbb{R}^n
\;\Big|\;
1+Y_i>0 \ \text{para todo } i
\Bigr\}.
\]
Para cada $Y\in C_{\mathrm{glob}}$, el vector de factores multiplicativos asociado viene dado por
\[
a(Y):=(1+Y_1,\dots,1+Y_n)\in\mathbb{R}_{>0}^n.
\]
\end{definicion}

\begin{definicion}[Operador informacional global]\label{def:Tglobal}
Sea $Y\in C_{\mathrm{glob}}$. El \emph{operador informacional global} asociado a $Y$ es la aplicación
\[
T_Y:\Delta_n^\circ\to\Delta_n^\circ,
\qquad
T_Y(X)
:=
\frac{X\circ a(Y)}{\sum_{j=1}^n X_j a_j(Y)},
\]
donde $\circ$ denota el producto coordenada a coordenada. En forma explícita,
\[
T_Y(X)_i
=
\frac{X_i(1+Y_i)}{\sum_{j=1}^n X_j(1+Y_j)},
\qquad i=1,\dots,n.
\]
\end{definicion}

\begin{observacion}\label{obs:Tglobal-frontera}
La fórmula de $T_Y$ tiene sentido también para $X\in\Delta_n$, no sólo para $X\in\Delta_n^\circ$. En efecto, si $Y\in C_{\mathrm{glob}}$, entonces $1+Y_i>0$ para todo $i$, y por tanto
\[
\sum_{j=1}^n X_j(1+Y_j)>0
\]
para todo $X\in\Delta_n$. Así, $T_Y(X)$ está bien definido en todo $\Delta_n$. No obstante, la restricción a $\Delta_n^\circ$ será la natural cuando se estudie la estructura de composición e inversión de estos operadores.
\end{observacion}
\begin{proposicion}[Preservación del interior del símplex]\label{prop:interiorGlobal}
Sea $Y\in C_{\mathrm{glob}}$. Entonces, para todo $X\in\Delta_n^\circ$,
\[
T_Y(X)\in\Delta_n^\circ.
\]
\end{proposicion}

\begin{proof}
Si $X_i>0$ y $a_i(Y)=1+Y_i>0$ para todo $i$, entonces
\[
X_i a_i(Y)>0,\qquad i=1,\dots,n.
\]
Además,
\[
\sum_{i=1}^n T_Y(X)_i
=
\frac{\sum_{i=1}^n X_i a_i(Y)}{\sum_{j=1}^n X_j a_j(Y)}
=
1.
\]
Por tanto, todas las coordenadas de $T_Y(X)$ son estrictamente positivas y su suma es $1$, es decir,
\[
T_Y(X)\in\Delta_n^\circ.
\]
\end{proof}

\begin{observacion}[Relación con el operador local del sección 2]\label{obs:localVsGlobal}
Sea $X\in\Delta_n^\circ$ y sea $Y\in C_{\mathrm{glob}}$. Si además
\[
\sum_{i=1}^n X_i(1+Y_i)=1,
\]
entonces
\[
T_Y(X)_i=X_i(1+Y_i),\qquad i=1,\dots,n,
\]
y por tanto
\[
T_Y(X)=T_Y^{\mathrm{loc}}(X).
\]
En este sentido, el operador global extiende el mecanismo local del sección 2 al incorporar una renormalización dependiente del estado base.
\end{observacion}

\begin{definicion}[Familia global de operadores]\label{def:G}
Se define
\[
\mathcal G
:=
\bigl\{
T_Y:\Delta_n^\circ\to\Delta_n^\circ
\;\big|\;
Y\in C_{\mathrm{glob}}
\bigr\},
\]
dotada de la composición usual de aplicaciones.
\end{definicion}

\begin{proposicion}[Ley de composición]\label{prop:composicionGlobal}
Sean $Y_1,Y_2\in C_{\mathrm{glob}}$. Entonces existe un único $Y_{\mathrm{tot}}\in C_{\mathrm{glob}}$ tal que
\[
T_{Y_2}\circ T_{Y_1}=T_{Y_{\mathrm{tot}}},
\]
y está determinado por la condición
\[
1+Y_{\mathrm{tot}}
=
(1+Y_1)\circ(1+Y_2).
\]
En coordenadas,
\[
1+(Y_{\mathrm{tot}})_i=(1+(Y_1)_i)(1+(Y_2)_i),
\qquad i=1,\dots,n.
\]
\end{proposicion}
\begin{proof}
Sean $Y^{(1)},Y^{(2)}\in C_{\mathrm{glob}}$. Para $r=1,2$, recordemos que el vector de factores multiplicativos esta dado por:
\[
a^{(r)}:=(a_1^{(r)},\dots,a_n^{(r)}),\qquad a_i^{(r)}:=1+Y_i^{(r)}.
\]
Como $Y^{(r)}\in C_{\mathrm{glob}}$, se tiene
\[
a_i^{(r)}>0\qquad \text{para todo } i=1,\dots,n.
\]

Sea ahora $X\in\Delta_n^\circ$. Entonces, por definición,
\[
T_{Y^{(1)}}(X)_i=\frac{X_i a_i^{(1)}}{\sum_{j=1}^n X_j a_j^{(1)}},\qquad i=1,\dots,n.
\]
Aplicando después $T_{Y^{(2)}}$, para cada $i=1,\dots,n$ se obtiene
\[
T_{Y^{(2)}}(T_{Y^{(1)}}(X))_i
=
\frac{T_{Y^{(1)}}(X)_i\,a_i^{(2)}}{\sum_{j=1}^n T_{Y^{(1)}}(X)_j\,a_j^{(2)}}.
\]
Sustituyendo la expresión de $T_{Y^{(1)}}(X)$,
\[
T_{Y^{(2)}}(T_{Y^{(1)}}(X))_i
=
\frac{
\displaystyle
\frac{X_i a_i^{(1)}}{\sum_{k=1}^n X_k a_k^{(1)}}\,a_i^{(2)}
}{
\displaystyle
\sum_{j=1}^n
\frac{X_j a_j^{(1)}}{\sum_{k=1}^n X_k a_k^{(1)}}\,a_j^{(2)}
}.
\]
El factor común $\sum_{k=1}^n X_k a_k^{(1)}$ aparece en numerador y denominador, por lo que se cancela. Así,
\[
T_{Y^{(2)}}(T_{Y^{(1)}}(X))_i
=
\frac{X_i a_i^{(1)} a_i^{(2)}}{\sum_{j=1}^n X_j a_j^{(1)} a_j^{(2)}}.
\]

Definimos ahora
\[
a^{(\mathrm{tot})}:=(a_1^{(1)}a_1^{(2)},\dots,a_n^{(1)}a_n^{(2)})\in \mathbb R_{>0}^n.
\]
Como $a_i^{(1)}>0$ y $a_i^{(2)}>0$ para todo $i$, se tiene
\[
a_i^{(\mathrm{tot})}>0\qquad \text{para todo } i.
\]
Definimos $Y^{(\mathrm{tot})}\in\mathbb R^n$ por
\[
1+Y_i^{(\mathrm{tot})}:=a_i^{(\mathrm{tot})},\qquad i=1,\dots,n.
\]
Entonces $Y^{(\mathrm{tot})}\in C_{\mathrm{glob}}$ y, para cada $i=1,\dots,n$,
\[
T_{Y^{(\mathrm{tot})}}(X)_i
=
\frac{X_i a_i^{(\mathrm{tot})}}{\sum_{j=1}^n X_j a_j^{(\mathrm{tot})}}
=
\frac{X_i a_i^{(1)} a_i^{(2)}}{\sum_{j=1}^n X_j a_j^{(1)} a_j^{(2)}}.
\]
Por consiguiente,
\[
T_{Y^{(2)}}(T_{Y^{(1)}}(X))_i=T_{Y^{(\mathrm{tot})}}(X)_i
\qquad \text{para todo } i=1,\dots,n,
\]
y por tanto
\[
T_{Y^{(2)}}\circ T_{Y^{(1)}}=T_{Y^{(\mathrm{tot})}}.
\]
\end{proof}

\begin{corolario}[Estructura de grupo abeliano]
El conjunto $\mathcal G$, dotado de la composición de aplicaciones, es un grupo abeliano.
\end{corolario}

\begin{proof}
Por la Proposición 3.2, si $Y^{(1)},Y^{(2)}\in C_{\mathrm{glob}}$, entonces la composición
\[
T_{Y^{(2)}}\circ T_{Y^{(1)}}
\]
vuelve a ser de la forma $T_{Y^{(\mathrm{tot})}}$, donde
\[
1+Y_i^{(\mathrm{tot})}=(1+Y_i^{(1)})(1+Y_i^{(2)}),
\qquad i=1,\dots,n.
\]
Por tanto, la composición en $\mathcal G$ corresponde al producto coordenado de los factores multiplicativos positivos.

Como el producto coordenado en $(0,\infty)^n$ es asociativo y conmutativo, la composición en $\mathcal G$ hereda estas propiedades.

El elemento neutro corresponde al vector
\[
a=(1,\dots,1),
\]
es decir, a $Y=0$, y en ese caso
\[
T_0(X)=X
\qquad \text{para todo } X\in\Delta_n^\circ.
\]

Ahora sea $Y\in C_{\mathrm{glob}}$ y sea
\[
a(Y)=(1+Y_1,\dots,1+Y_n)\in(0,\infty)^n.
\]
Definimos
\[
a(Y)^{-1}
:=
\left(
\frac{1}{1+Y_1},\dots,\frac{1}{1+Y_n}
\right)\in(0,\infty)^n.
\]
Sea $Z\in\mathbb R^n$ dado por
\[
1+Z_i:=\frac{1}{1+Y_i},
\qquad i=1,\dots,n.
\]
Entonces $1+Z_i>0$ para todo $i$, de modo que $Z\in C_{\mathrm{glob}}$. Además,
\[
(1+Y_i)(1+Z_i)=1,
\qquad i=1,\dots,n.
\]
Por la ley de composición,
\[
T_Z\circ T_Y=T_0.
\]
Análogamente,
\[
T_Y\circ T_Z=T_0.
\]
Así, todo elemento de $\mathcal G$ es invertible.

En consecuencia, $\mathcal G$ es un grupo abeliano.
\end{proof}

\begin{observacion}[Coordenadas log-multiplicativas]\label{obs:coordenadasLog}
Para $Y\in C_{\mathrm{glob}}$ se definen las coordenadas
\[
\theta(Y) := (\theta_1,\dots,\theta_n), \qquad \theta_i := \log(1+Y_i), \quad i=1,\dots,n.
\]
Entonces
\[
1+Y_i=e^{\theta_i},
\]
y el operador global puede escribirse como
\[
T_\theta(X)
=
\frac{X\circ e^\theta}{\sum_{j=1}^n X_j e^{\theta_j}},
\qquad
e^\theta:=(e^{\theta_1},\dots,e^{\theta_n}).
\]
Bajo esta parametrización, la ley de composición del corolario anterior se traduce en
\[
\theta_{\mathrm{tot}}=\theta^{(1)}+\theta^{(2)}.
\]
En este sentido, la estructura multiplicativa de los factores positivos se convierte en una estructura aditiva en coordenadas logarítmicas la cual es compatible con el uso de coordenadas relativas en análisis composicional (\cite{Aitchison1986}).
\end{observacion}

\section{Sección D — Divergencia KL, métrica de Fisher y dinámica inducida}

En esta sección  se introduce el funcional informacional
\[
E_P(X):=D_{\mathrm{KL}}(P\|X),
\]
se calcula su gradiente respecto de la métrica de Fisher en el interior del símplex y se obtiene la dinámica continua asociada. Finalmente, se muestra que esta dinámica es compatible con la familia de operadores globales introducida en la sección 3 .

El resultado principal de la sección  será el Teorema~\ref{teo:equivalenciaCD}, donde se establece la compatibilidad estructural entre la dinámica Fisher--KL y la familia de operadores globales de la sección 3.

\begin{definicion}[Divergencia de Kullback--Leibler discreta]\label{def:KL}
Sea $P\in\Delta_n^\circ$ fija. Para $X\in\Delta_n$ se define
\[
D_{\mathrm{KL}}(P\|X)
=
\begin{cases}
\displaystyle \sum_{i=1}^n P_i\log\!\left(\frac{P_i}{X_i}\right),
& \text{si } X_i>0 \text{ para todo } i,\\[1em]
+\infty,
& \text{si existe } i \text{ tal que } P_i>0 \text{ y } X_i=0.
\end{cases}
\]

Se define además el funcional
\[
E_P:\Delta_n^\circ\to\mathbb R,
\qquad
E_P(X):=D_{\mathrm{KL}}(P\|X).
\]
\end{definicion}

\begin{observacion}[Dominio natural de $E_P$]\label{obs:KLsoporte}
Como $P\in\Delta_n^\circ$, se tiene $P_i>0$ para todo $i$. Por tanto,
\[
D_{\mathrm{KL}}(P\|X)<\infty
\quad\Longleftrightarrow\quad
X\in\Delta_n^\circ.
\]
En consecuencia, el dominio natural de $E_P$ es el interior del símplex.
\end{observacion}

\begin{definicion}[Espacio tangente del símplex]\label{def:tangente}
Sea $X\in\Delta_n^\circ$. El \emph{espacio tangente en $X$} se define como el conjunto de velocidades de curvas suaves del símplex que pasan por $X$, es decir,
\[
T_X\Delta_n^\circ
:=
\left\{
u\in\mathbb R^n
\;\middle|\;
\exists\,\varepsilon>0,\ \exists\,\gamma:(-\varepsilon,\varepsilon)\to\Delta_n^\circ
\ \text{de clase } C^1,\ 
\gamma(0)=X,\ 
\gamma'(0)=u
\right\}.
\]
\end{definicion}

\begin{observacion}
Para todo $X\in\Delta_n^\circ$,
\[
T_X\Delta_n^\circ
=
\left\{
u=(u_1,\dots,u_n)\in\mathbb R^n
\;\middle|\;
\sum_{i=1}^n u_i=0
\right\}.
\]
\end{observacion}

\begin{definicion}[Métrica de Fisher discreta]\label{def:Fisher}
Para $X\in\Delta_n^\circ$, la \emph{métrica de Fisher} en $T_X\Delta_n$ es la forma bilineal
\[
g_X(u,v)
:=
\sum_{i=1}^n \frac{u_i v_i}{X_i},
\qquad
u,v\in T_X\Delta_n.
\]
Esta es la realización discreta de la métrica de Fisher--Rao en la variedad de distribuciones finitas (\cite{Chentsov1982,AmariNagaoka2000,Amari2016}).
\end{definicion}

\begin{definicion}[Gradiente riemanniano respecto de la métrica de Fisher]\label{def:gradienteFisher}
Sea $E:\Delta_n^\circ\to\mathbb R$ una función suave. El \emph{gradiente riemanniano} de $E$ respecto de $g$ es el campo
\[
\mathrm{grad}_gE:\Delta_n^\circ\to T\Delta_n,
\qquad
X\mapsto \mathrm{grad}_gE(X)\in T_X\Delta_n,
\]
caracterizado por
\[
g_X(\mathrm{grad}_gE(X),u)=\mathrm dE(X)[u],
\qquad
\forall\,u\in T_X\Delta_n.
\]
\end{definicion}
\begin{proposicion}[Gradiente Fisher de la divergencia KL]\label{prop:gradKL}
Sea $P\in\Delta_n^\circ$ fija y sea
\[
E_P(X)=D_{\mathrm{KL}}(P\|X)=\sum_{i=1}^n P_i \log\!\left(\frac{P_i}{X_i}\right),
\qquad X\in\Delta_n^\circ.
\]
Entonces
\[
\operatorname{grad}_g E_P(X)=X-P,
\qquad X\in\Delta_n^\circ.
\]
\end{proposicion}

\begin{proof}
Fijemos $X\in\Delta_n^\circ$ y sea $u\in T_X\Delta_n^\circ$. Entonces
\[
dE_P(X)[u]
=
\left.\frac{d}{dt}E_P(X+tu)\right|_{t=0}.
\]
Como
\[
E_P(X+tu)
=
\sum_{i=1}^n P_i \log\!\left(\frac{P_i}{X_i+tu_i}\right),
\]
se obtiene
\[
\frac{d}{dt}E_P(X+tu)
=
-\sum_{i=1}^n P_i \frac{u_i}{X_i+tu_i}.
\]
Evaluando en $t=0$,
\[
dE_P(X)[u]
=
-\sum_{i=1}^n \frac{P_i}{X_i}u_i.
\]

Por otra parte, por definición de la métrica de Fisher,
\[
g_X(v,u)=\sum_{i=1}^n \frac{v_i u_i}{X_i}.
\]
Tomando $v=X-P$, se tiene
\[
g_X(X-P,u)
=
\sum_{i=1}^n \frac{(X_i-P_i)u_i}{X_i}
=
\sum_{i=1}^n u_i - \sum_{i=1}^n \frac{P_i}{X_i}u_i.
\]
Como $u\in T_X\Delta_n^\circ$, se cumple $\sum_{i=1}^n u_i=0$, y por tanto
\[
g_X(X-P,u)
=
-\sum_{i=1}^n \frac{P_i}{X_i}u_i
=
dE_P(X)[u].
\]
Por la definición del gradiente riemanniano, se concluye que
\[
\operatorname{grad}_g E_P(X)=X-P.
\]
\end{proof}

\begin{definicion}[Flujo de gradiente Fisher--KL]\label{def:flujoFisherKL}
El \emph{flujo de gradiente} asociado a $E_P$ es la ecuación diferencial
\[
\dot X(t)=-\mathrm{grad}_gE_P(X(t)),
\qquad
X(t)\in\Delta_n^\circ.
\]
Por la Proposición~\ref{prop:gradKL}, esta ecuación toma la forma
\[
\dot X(t)=P-X(t).
\]
\end{definicion}

\begin{proposicion}[Solución explícita del flujo]\label{prop:solucionExplicita}
Sea $X_0\in\Delta_n^\circ$. La solución del problema
\[
\dot X(t)=P-X(t),
\qquad
X(0)=X_0,
\]
está dada por
\[
X(t)=e^{-t}X_0+(1-e^{-t})P,
\qquad t\ge 0.
\]
En particular, $X(t)\in\Delta_n^\circ$ para todo $t\ge 0$ y
\[
X(t)\longrightarrow P
\qquad (t\to\infty).
\]
\end{proposicion}

\begin{proof}
La ecuación
\[
\dot X(t)=P-X(t),\qquad X(0)=X_0,
\]
se desacopla coordenada a coordenada. En efecto, para cada $i=1,\dots,n$,
\[
\dot X_i(t)=P_i-X_i(t),\qquad X_i(0)=(X_0)_i.
\]
La solución de esta ecuación lineal escalar es
\[
X_i(t)=e^{-t}(X_0)_i+(1-e^{-t})P_i,
\qquad i=1,\dots,n.
\]
Por tanto,
\[
X(t)=e^{-t}X_0+(1-e^{-t})P.
\]

Como $X_0\in\Delta_n^\circ$ y $P\in\Delta_n^\circ$, se tiene
\[
(X_0)_i>0,\qquad P_i>0,\qquad i=1,\dots,n.
\]
Además, para todo $t\ge 0$,
\[
e^{-t}>0,
\qquad
1-e^{-t}\ge 0.
\]
De aquí se sigue que, para cada $i=1,\dots,n$,
\[
X_i(t)=e^{-t}(X_0)_i+(1-e^{-t})P_i>0.
\]
Por otra parte,
\[
\sum_{i=1}^n X_i(t)
=
e^{-t}\sum_{i=1}^n (X_0)_i
+
(1-e^{-t})\sum_{i=1}^n P_i
=
e^{-t}+(1-e^{-t})
=
1.
\]
Por tanto,
\[
X(t)\in\Delta_n^\circ
\qquad \text{para todo } t\ge 0.
\]

Finalmente,
\[
X(t)-P
=
e^{-t}X_0+(1-e^{-t})P-P
=
e^{-t}(X_0-P).
\]
Como $e^{-t}\to 0$ cuando $t\to\infty$, se concluye que
\[
X(t)\to P
\qquad (t\to\infty).
\]
\end{proof}

\begin{proposicion}[Monotonicidad de $E_P$ a lo largo del flujo]\label{prop:Hteorema}
Sea $X(\cdot)$ una solución del flujo Fisher--KL. Entonces la función
\[
t\longmapsto E_P(X(t))
\]
es decreciente y satisface
\[
\frac{d}{dt}E_P(X(t))
=
-\,g_{X(t)}\!\bigl(\mathrm{grad}_gE_P(X(t)),\mathrm{grad}_gE_P(X(t))\bigr)
\le 0.
\]
La igualdad sólo ocurre si $X(t)=P$.
\end{proposicion}

\begin{proof}
Por la regla de la cadena,
\[
\frac{d}{dt}E_P(X(t))
=
\mathrm dE_P(X(t))[\dot X(t)].
\]
Como
\[
\dot X(t)=-\mathrm{grad}_gE_P(X(t)),
\]
y por definición del gradiente riemanniano,
\[
\mathrm dE_P(X(t))[\dot X(t)]
=
-\,g_{X(t)}\!\bigl(\mathrm{grad}_gE_P(X(t)),\mathrm{grad}_gE_P(X(t))\bigr).
\]
La no positividad es inmediata porque $g$ es definida positiva en cada espacio tangente. Además,
\[
g_{X(t)}\!\bigl(\mathrm{grad}_gE_P(X(t)),\mathrm{grad}_gE_P(X(t))\bigr)=0
\]
si y sólo si
\[
\mathrm{grad}_gE_P(X(t))=0
\quad\Longleftrightarrow\quad
X(t)-P=0,
\]
por la Proposición~\ref{prop:gradKL}. Por tanto, la igualdad sólo ocurre en el equilibrio $X(t)=P$.
\end{proof}

\begin{proposicion}[Mínimo global de $E_P$]\label{prop:minimoGlobalKL}
Para todo $X\in\Delta_n^\circ$ se cumple
\[
E_P(X)\ge 0,
\]
y además
\[
E_P(X)=0
\quad\Longleftrightarrow\quad
X=P.
\]
En particular, $P$ es el único minimizador global de $E_P$ en $\Delta_n^\circ$.
\end{proposicion}

\begin{proof}
La no negatividad y el criterio de igualdad corresponden a la desigualdad de Gibbs para la divergencia de Kullback--Leibler (\cite{Csiszar1975,CoverThomas1991}). Aplicando ese resultado a
\[
E_P(X)=D_{\mathrm{KL}}(P\|X),
\]
se obtiene inmediatamente
\[
E_P(X)\ge 0
\quad\text{y}\quad
E_P(X)=0 \Longleftrightarrow X=P.
\]
De aquí se sigue que $P$ es el único minimizador global.
\end{proof}

\begin{teorema}[Equivalencia estructural Sección C $\leftrightarrow$ Sección D]
\label{teo:equivalenciaCD}
Sea $P\in\Delta_n^\circ$ fijo. Considérese el operador informacional global $T_Y$
del sección ~C (Def.~\ref{def:Tglobal}), definido para $Y\in C_{\mathrm{glob}}$
(Def.~\ref{def:Cglob}) por
\[
T_Y(X)=\frac{X\circ(1+Y)}{\langle X,1+Y\rangle}.
\]
Defínase, para cada $X\in\Delta_n^\circ$, el campo de corrección óptima
\[
Y^*(X,P)\in\mathbb{R}^n,
\qquad
Y^*_i(X,P):=\frac{P_i}{X_i}-1,\quad i=1,\dots,n.
\]
Entonces se verifican las siguientes afirmaciones.

\smallskip
\noindent\textbf{(i) Nivel algebraico (corrección multiplicativa).}
Para todo $X\in\Delta_n^\circ$ se tiene $Y^*(X,P)\in C(X)$ (Def.~\ref{def:CX}) y
\[
T_{Y^*(X,P)}(X)=P.
\]

\smallskip
\noindent\textbf{(ii) Nivel discreto (dinámica inducida).}
Fijado $\alpha\in(0,1]$, la iteración
\[
X_{k+1}:=T_{\alpha\,Y^*(X_k,P)}(X_k),\qquad X_0\in\Delta_n^\circ,
\]
se reduce a la interpolación convexa
\[
X_{k+1}=(1-\alpha)X_k+\alpha P,
\]
y por tanto permanece en $\Delta_n^\circ$ (Prop.~\ref{prop:interiorGlobal}).

\smallskip
\noindent\textbf{(iii) Nivel continuo (límite infinitesimal y geometría Fisher).}
El límite de primer orden $\alpha\to 0$ de la regla
$T_{\alpha Y}(X)$ conduce a la ecuación diferencial
\[
\dot X = X\circ Y - X\,\langle X,Y\rangle.
\]
Aplicada al campo óptimo $Y=Y^*(X,P)$, esta ecuación se reduce a
\[
\dot X = P-X,
\]
que coincide con el flujo de gradiente Fisher--KL del funcional $E_P$
(Def.~\ref{def:flujoFisherKL}), es decir,
\[
\dot X=-\mathrm{grad}_g E_P(X)
\qquad
(\text{Prop.~\ref{prop:gradKL}}).
\]

\end{teorema}

\begin{proof}
\noindent\textbf{(i)} Sea $X\in\Delta_n^\circ$. Como $X_i>0$ y $P_i>0$ para todo $i$,
se tiene $1+Y^*_i(X,P)=P_i/X_i>0$, luego $Y^*(X,P)\in C_{\mathrm{glob}}$
(Def.~\ref{def:Cglob}). Además,
\[
X_i\bigl(1+Y^*_i(X,P)\bigr)=P_i\ge 0,
\qquad
\sum_{i=1}^n X_i\bigl(1+Y^*_i(X,P)\bigr)=\sum_{i=1}^n P_i=1,
\]
luego $Y^*(X,P)\in C(X)$ (Def.~\ref{def:CX}). Sustituyendo en el operador global,
\[
T_{Y^*(X,P)}(X)_i
=
\frac{X_i(P_i/X_i)}{\sum_{j=1}^n X_j(P_j/X_j)}
=
\frac{P_i}{\sum_{j=1}^n P_j}
=
P_i,
\]
y por tanto $T_{Y^*(X,P)}(X)=P$.

\smallskip
\noindent\textbf{(ii)} Sea $\alpha\in(0,1]$ y $X\in\Delta_n^\circ$. Entonces
\[
1+\alpha Y^*_i(X,P)
=
(1-\alpha)+\alpha\,\frac{P_i}{X_i},
\]
y al sustituir en $T_{\alpha Y^*}(X)$ se obtiene, coordenada a coordenada,
\[
T_{\alpha Y^*(X,P)}(X)_i
=
\frac{(1-\alpha)X_i+\alpha P_i}{(1-\alpha)\sum_{j}X_j+\alpha\sum_{j}P_j}
=
(1-\alpha)X_i+\alpha P_i,
\]
pues $\sum_j X_j=\sum_j P_j=1$. Iterando, $X_{k+1}=(1-\alpha)X_k+\alpha P$, y por
ser combinación convexa de puntos de $\Delta_n^\circ$ se mantiene
$X_k\in\Delta_n^\circ$ (Prop.~\ref{prop:interiorGlobal}).

\smallskip
\noindent\textbf{(iii)} Partimos de
\[
T_{\alpha Y}(X)
=
\frac{X+\alpha(X\circ Y)}{1+\alpha\langle X,Y\rangle}.
\]
Usando la expansión $\frac{1}{1+\alpha a}=1-\alpha a+o(\alpha)$,
\[
T_{\alpha Y}(X)
=
X+\alpha\bigl(X\circ Y-X\langle X,Y\rangle\bigr)+o(\alpha),
\]
lo que conduce a la ecuación diferencial anunciada.

Tomando $Y=Y^*(X,P)$,
\[
X\circ Y^*(X,P)=P-X,
\qquad
\langle X,Y^*(X,P)\rangle=\sum_i P_i-\sum_i X_i=0,
\]
y sustituyendo se obtiene $\dot X=P-X$. Por la
Proposición~\ref{prop:gradKL}, $\mathrm{grad}_g E_P(X)=X-P$, de modo que
\[
\dot X=P-X=-\mathrm{grad}_g E_P(X).
\]
Esto establece la equivalencia estructural entre los tres niveles.
\end{proof}


% ============================================================
\section{Sección E\textsubscript{0} — Semántica probabilística: Bayes y Markov como casos particulares}
% ============================================================

Esta sección dota al marco de Control Informacional de la semántica necesaria para que los operadores de las Secciones B--D puedan interpretarse y utilizarse como la actualización bayesiana clásica y el operador de Markov clásico, respectivamente. El objetivo no es ampliar el marco sino identificar, dentro de él, los casos particulares que corresponden a esos dos paradigmas, de manera que los problemas clásicos de Bayes y Markov puedan plantearse y resolverse en el lenguaje de CI.

\subsection{Semántica bayesiana}

La fundamentación bayesiana se construye en cuatro pasos, todos dentro del marco ya establecido. Primero se define el espacio conjunto hipótesis-observaciones como un espacio de probabilidad en el sentido de la Sección~A. Segundo, se identifica la verosimilitud como la corrección global inducida por la probabilidad condicional sobre ese espacio. Tercero, se prueba que $T_{Y^\ell}$ coincide con el posterior de Bayes. Cuarto, se establece la unicidad: $T_{Y^\ell}$ es el único operador en $\mathcal{G}$ que es coherente con la probabilidad condicional del espacio conjunto.

\begin{definicion}[Espacio conjunto hipótesis-observaciones]\label{def:espacio_conjunto}
Sea $(\Omega,\mathcal{F},P)$ un espacio de probabilidad (Definición~\ref{def:espacio_prob}) y sea
\[
\Pi_H = \{H_1,\dots,H_n\} \subset \mathcal{F}
\]
una partición medible finita no degenerada de $\Omega$, cuyos elementos se interpretan como \emph{hipótesis alternativas}. Sea $\Pi_E = \{E_1,\dots,E_m\} \subset \mathcal{F}$ una segunda partición medible finita no degenerada de $\Omega$, cuyos elementos se interpretan como \emph{observaciones posibles}.

El \emph{espacio conjunto} es el mismo espacio $(\Omega,\mathcal{F},P)$, equipado con las dos particiones simultáneamente. El estado informacional previo (\emph{prior}) es el vector
\[
X_i := P(H_i), \qquad i=1,\dots,n,
\]
que por la Proposición~\ref{prop:estado_simplex} pertenece a $\Delta_n^\circ$.
\end{definicion}

\begin{observacion}[Por qué una sola medida $P$]\label{obs:una_medida}
Las dos particiones $\Pi_H$ y $\Pi_E$ viven sobre el mismo espacio $(\Omega,\mathcal{F},P)$. Esto permite definir todas las probabilidades conjuntas y condicionales necesarias directamente desde $P$, sin recurrir a estructuras adicionales. En particular, $P(H_i \cap E_j)$, $P(H_i \mid E_j)$ y $P(E_j \mid H_i)$ son todos objetos bien definidos dentro de la Definición~\ref{def:espacio_prob}.
\end{observacion}

\begin{definicion}[Verosimilitud como corrección global]\label{def:verosimilitud}
Sea $j\in\{1,\dots,m\}$ un índice de observación con $P(E_j)>0$. La \emph{verosimilitud de la observación $E_j$ bajo la hipótesis $H_i$} es
\[
\ell_i^{(j)} := P(E_j \mid H_i) = \frac{P(H_i \cap E_j)}{P(H_i)}, \qquad i=1,\dots,n.
\]
Como $P(H_i)>0$ (partición no degenerada) y $P(H_i\cap E_j)\geq 0$, se tiene $\ell_i^{(j)}\geq 0$. La condición $\ell_i^{(j)}>0$ para todo $i$ equivale a que la observación $E_j$ sea compatible con cada hipótesis, es decir, $P(H_i\cap E_j)>0$ para todo $i$.

Cuando esta condición se satisface, el vector $\ell^{(j)} = (\ell_1^{(j)},\dots,\ell_n^{(j)})\in\mathbb{R}_{>0}^n$ define una corrección global $Y^{\ell^{(j)}}\in C_{\mathrm{glob}}$ (Definición~\ref{def:Cglob}) mediante
\[
1 + Y^{\ell^{(j)}}_i := \ell^{(j)}_i = P(E_j\mid H_i).
\]
\end{definicion}

\begin{proposicion}[El operador bayesiano como caso particular de $T_Y$]\label{prop:bayes_como_TY}
Sea $X\in\Delta_n^\circ$ el prior con $X_i=P(H_i)$, y sea $\ell^{(j)}\in\mathbb{R}_{>0}^n$ la verosimilitud de $E_j$. Entonces
\[
T_{Y^{\ell^{(j)}}}(X)_i
= \frac{X_i\,\ell^{(j)}_i}{\displaystyle\sum_{k=1}^n X_k\,\ell^{(j)}_k}
= P(H_i \mid E_j),
\qquad i=1,\dots,n.
\]
Es decir, $T_{Y^{\ell^{(j)}}}(X)$ es el estado informacional posterior dado que se observó $E_j$.
\end{proposicion}

\begin{proof}
Por la regla de Bayes y la ley de probabilidad total,
\[
P(H_i\mid E_j)
= \frac{P(E_j\mid H_i)\,P(H_i)}{P(E_j)}
= \frac{\ell^{(j)}_i\,X_i}{\displaystyle\sum_{k=1}^n P(E_j\mid H_k)\,P(H_k)}
= \frac{\ell^{(j)}_i\,X_i}{\displaystyle\sum_{k=1}^n \ell^{(j)}_k\,X_k}.
\]
Por la Definición~\ref{def:Tglobal} con factores $a_i = 1+Y^{\ell^{(j)}}_i = \ell^{(j)}_i$,
\[
T_{Y^{\ell^{(j)}}}(X)_i = \frac{X_i\,\ell^{(j)}_i}{\sum_{k=1}^n X_k\,\ell^{(j)}_k}.
\]
Ambas expresiones coinciden.
\end{proof}

\begin{proposicion}[Unicidad: $T_{Y^{\ell^{(j)}}}$ es el único operador coherente con $P$]\label{prop:unicidad_bayes}
Sea $(\Omega,\mathcal{F},P)$ el espacio conjunto de la Definición~\ref{def:espacio_conjunto} con particiones $\Pi_H$ y $\Pi_E$. Sea $j\in\{1,\dots,m\}$ con $P(H_i\cap E_j)>0$ para todo $i$. Un operador $T\in\mathcal{G}$ satisface
\[
T(X)_i = P(H_i\mid E_j) \qquad \text{para todo } i=1,\dots,n
\]
si y sólo si $T = T_{Y^{\ell^{(j)}}}$, donde $Y^{\ell^{(j)}}$ es la corrección de la Definición~\ref{def:verosimilitud}.

En otras palabras, dentro de $\mathcal{G}$, la actualización bayesiana es la única que produce el posterior correcto determinado por $P$.
\end{proposicion}

\begin{proof}
La suficiencia es la Proposición~\ref{prop:bayes_como_TY}.

Para la necesidad: supóngase que $T_Y\in\mathcal{G}$ satisface $T_Y(X)_i = P(H_i\mid E_j)$ para todo $i$. Por la Definición~\ref{def:Tglobal},
\[
T_Y(X)_i = \frac{X_i(1+Y_i)}{\sum_{k=1}^n X_k(1+Y_k)}.
\]
Igualando con $P(H_i\mid E_j) = \ell^{(j)}_i X_i / \sum_k \ell^{(j)}_k X_k$, se obtiene para cada $i$:
\[
\frac{X_i(1+Y_i)}{\sum_{k} X_k(1+Y_k)} = \frac{X_i \ell^{(j)}_i}{\sum_{k} X_k \ell^{(j)}_k}.
\]
Como $X_i>0$, dividiendo numerador y denominador:
\[
\frac{1+Y_i}{\sum_{k} X_k(1+Y_k)} = \frac{\ell^{(j)}_i}{\sum_{k} X_k \ell^{(j)}_k}.
\]
Sea $c := \frac{\sum_k X_k(1+Y_k)}{\sum_k X_k\ell^{(j)}_k} > 0$. Entonces $1+Y_i = c\,\ell^{(j)}_i$ para todo $i$, es decir, $Y = Y^{c\ell^{(j)}}$. Por la invariancia de escala (Proposición~\ref{prop:composicionGlobal}), $T_{Y^{c\ell^{(j)}}} = T_{Y^{\ell^{(j)}}}$. Luego $T_Y = T_{Y^{\ell^{(j)}}}$.
\end{proof}

\begin{observacion}[Alcance de la unicidad]\label{obs:alcance_unicidad}
La Proposición~\ref{prop:unicidad_bayes} establece unicidad dentro de $\mathcal{G}$: dado el espacio conjunto $(\Omega,\mathcal{F},P)$, el posterior $P(H_i\mid E_j)$ está unívocamente determinado, y el único operador global que lo produce es $T_{Y^{\ell^{(j)}}}$. Lo que no establece es que $\mathcal{G}$ sea la única familia de operadores razonable en general; eso requeriría un argumento de coherencia normativa (tipo Cox o De Finetti) que queda fuera del alcance de este trabajo. Lo que sí se garantiza es que, dentro del marco de CI, la regla de Bayes no admite alternativas en $\mathcal{G}$.
\end{observacion}

\begin{corolario}[Actualización bayesiana secuencial como composición en $\mathcal{G}$]\label{cor:bayes_secuencial}
Sean $E_{j_1},\dots,E_{j_m}$ observaciones sucesivas con $P(H_i\cap E_{j_s})>0$ para todo $i,s$. Sea $X_0=X\in\Delta_n^\circ$ el prior inicial. Por la Proposición~\ref{prop:unicidad_bayes}, cada paso $s$ produce el único operador coherente con $P(\,\cdot\mid E_{j_s})$. El posterior tras las $m$ observaciones es
\[
X_m = T_{Y^{\ell^{(j_m)}}}\circ\cdots\circ T_{Y^{\ell^{(j_1)}}}(X_0),
\]
y por la Proposición~\ref{prop:composicionGlobal} este operador compuesto tiene factores
\[
1+(Y^{\mathrm{tot}})_i = \prod_{s=1}^m \ell^{(j_s)}_i.
\]
Si además las observaciones son condicionalmente independientes dadas las hipótesis, es decir,
\[
P(E_{j_1}\cap\cdots\cap E_{j_m}\mid H_i) = \prod_{s=1}^m P(E_{j_s}\mid H_i),
\]
entonces $Y^{\mathrm{tot}}$ coincide con la corrección de la observación conjunta $E_{j_1}\cap\cdots\cap E_{j_m}$, y la actualización secuencial es idéntica a la simultánea.

En coordenadas logarítmicas (Observación~\ref{obs:coordenadasLog}), la acumulación de evidencia es aditiva:
\[
\log(1+(Y^{\mathrm{tot}})_i) = \sum_{s=1}^m \log \ell^{(j_s)}_i.
\]
\end{corolario}

\begin{proof}
La forma del operador compuesto es la Proposición~\ref{prop:composicionGlobal}. La identificación del operador compuesto con la observación conjunta bajo independencia condicional se sigue de que $P(E_{j_1}\cap\cdots\cap E_{j_m}\mid H_i) = \prod_s\ell^{(j_s)}_i$, que es precisamente el factor del operador compuesto.
\end{proof}

\begin{proposicion}[Corrección óptima hacia el posterior]\label{prop:correccion_optima_bayes}
Sean $X,\pi\in\Delta_n^\circ$. El único elemento de $\mathcal{G}$ (salvo escala proyectiva) que lleva $X$ a $\pi$ en un paso es $T_{Y^*}$ con
\[
1+Y^*_i = c\,\frac{\pi_i}{X_i}, \qquad c>0.
\]
Si $\pi = P(H_i\mid E_j)$ es el posterior, entonces por la Proposición~\ref{prop:unicidad_bayes} este operador coincide con $T_{Y^{\ell^{(j)}}}$, y la corrección óptima satisface $\ell^{(j)}_i \propto \pi_i/X_i$.
\end{proposicion}

\begin{proof}
Por el Teorema~\ref{teo:equivalenciaCD}(i), el campo que lleva $X$ a $\pi$ satisface $1+Y^*_i = \pi_i/X_i$. La unicidad salvo escala se sigue de la estructura de grupo de $\mathcal{G}$ (Corolario de la Proposición~\ref{prop:composicionGlobal}). La identificación con $\ell^{(j)}$ se sigue de la Proposición~\ref{prop:unicidad_bayes}: el único operador de $\mathcal{G}$ que produce $P(H_i\mid E_j)$ desde $X$ tiene factores proporcionales a $P(H_i\mid E_j)/X_i = \ell^{(j)}_i/\sum_k\ell^{(j)}_k X_k$.
\end{proof}

\subsection{Puente estructural hacia Markov}

Antes de introducir la semántica markoviana conviene precisar por qué las cadenas de Markov entran en CI a través de la Sección~B y no de la Sección~C. La diferencia es estructural: la actualización bayesiana repondera coordenadas por factores fijos independientes del estado, mientras que un operador de Markov redistribuye masa entre coordenadas mediante una matriz, produciendo una corrección que depende del estado base.

\begin{definicion}[Conjunto de matrices estocásticas por filas]\label{def:Sn_markov}
Se define
\[
\mathcal{S}_n
:=
\left\{
K\in\mathbb{R}^{n\times n}
\;\middle|\;
K_{ij}\ge 0\ \text{para todo } i,j,
\quad
\sum_{j=1}^n K_{ij}=1\ \text{para todo } i
\right\}.
\]
Los elementos de $\mathcal{S}_n$ se llaman \emph{matrices estocásticas por filas} o \emph{matrices de transición}.
\end{definicion}

\begin{proposicion}[El símplex es invariante bajo $\mathcal{S}_n$]\label{prop:simplex_invariante}
Sea $K\in\mathcal{S}_n$ y sea $X\in\Delta_n$. Entonces $XK\in\Delta_n$. En particular, toda $K\in\mathcal{S}_n$ induce un operador bien definido
\[
L_K:\Delta_n\longrightarrow\Delta_n,\qquad L_K(X):=XK.
\]
\end{proposicion}

\begin{proof}
Para cada $j$,
\[
(XK)_j = \sum_{i=1}^n X_i K_{ij} \ge 0,
\]
pues $X_i\ge 0$ y $K_{ij}\ge 0$. Además,
\[
\sum_{j=1}^n (XK)_j
=\sum_{j=1}^n\sum_{i=1}^n X_iK_{ij}
=\sum_{i=1}^n X_i\underbrace{\sum_{j=1}^n K_{ij}}_{=1}
=\sum_{i=1}^n X_i=1.
\]
Luego $XK\in\Delta_n$.
\end{proof}

\begin{observacion}[Sobre la preservación del interior]\label{obs:interior_markov}
La Proposición~\ref{prop:simplex_invariante} garantiza $L_K(\Delta_n)\subset\Delta_n$, pero no $L_K(\Delta_n^\circ)\subset\Delta_n^\circ$ para una $K\in\mathcal{S}_n$ general. Una condición \emph{suficiente} para preservar el interior es $K_{ij}>0$ para todo $i,j$, pero no es necesaria: la matriz identidad, por ejemplo, preserva $\Delta_n^\circ$ y tiene ceros fuera de la diagonal. En lo que sigue se trabajará con $X_t\in\Delta_n$ salvo que se indique explícitamente una hipótesis de positividad. La corrección local $Y^K(X)$ requiere $X\in\Delta_n^\circ$ para que los cocientes $(XK)_i/X_i$ estén definidos, por lo que en esa proposición se impondrá dicha hipótesis.
\end{observacion}

\begin{proposicion}[Markov entra por la Sección B, no por la Sección C]\label{prop:markov_es_local}
Sea $K\in\mathcal{S}_n$ y sea $X\in\Delta_n^\circ$. El paso markoviano $L_K(X)=XK$ puede representarse como una corrección local en el sentido de la Sección~B mediante
\[
Y^K_i(X) := \frac{(XK)_i}{X_i}-1, \qquad i=1,\dots,n.
\]
Se tiene $Y^K(X)\in C(X)$ y $T^{\mathrm{loc}}_{Y^K(X)}(X)=XK$.

Por el contrario, salvo el caso trivial $K=I$, el operador $L_K$ no puede representarse como un operador global $T_Y\in\mathcal{G}$ con $Y$ independiente de $X$.
\end{proposicion}

\begin{proof}
Para cada $i$: $X_i(1+Y^K_i(X))=(XK)_i\ge 0$, y $\sum_i X_i(1+Y^K_i(X))=\sum_i(XK)_i=1$. Luego $Y^K(X)\in C(X)$ y $T^{\mathrm{loc}}_{Y^K(X)}(X)_i=(XK)_i$ por la Definición~\ref{def:Tloc}.

Para la segunda parte: supóngase que existe $a\in\mathbb{R}_{>0}^n$ tal que $T_a(X)=XK$ para todo $X\in\Delta_n^\circ$. Si $a_i=c$ para todo $i$, entonces $T_a(X)_i=X_i$ para todo $X$, que coincide con $XK$ sólo si $K=I$. Si $a$ no es proporcional a $(1,\dots,1)$, entonces $T_a(X)$ depende de $X$ de forma no lineal a través del denominador $\sum_k X_k a_k$, mientras que $(XK)_i=\sum_j X_j K_{ji}$ es lineal en $X$; estas dos funciones no pueden coincidir para todo $X\in\Delta_n^\circ$ salvo en el caso degenerado anterior. En consecuencia, para $K\neq I$ no existe $Y\in C_{\mathrm{glob}}$ independiente de $X$ tal que $T_Y(X)=XK$ para todo $X$.
\end{proof}

\begin{observacion}[La ruta conceptual de cada paradigma]\label{obs:rutas_paradigmas}
El puente queda así fijado:
\[
\underbrace{A: X\in\Delta_n}_{\text{dominio}}
\;\longrightarrow\;
\underbrace{B: T^{\mathrm{loc}}_{Y^K(X)}}_{\text{Markov}}
\qquad\text{y}\qquad
\underbrace{A: X\in\Delta_n^\circ}_{\text{dominio}}
\;\longrightarrow\;
\underbrace{C: T_{Y^\ell}\in\mathcal{G}}_{\text{Bayes}}.
\]
La Sección~D interviene en ambos casos: para Bayes, el funcional $E_P(X)=D_{\mathrm{KL}}(P\|X)$ mide la distancia al posterior y el flujo $\dot X=P-X$ describe la convergencia continua; para Markov, el mismo funcional con $P=\pi$ estacionaria mide el desfase hacia el equilibrio y la Proposición~\ref{prop:disipacion_KL_markov} garantiza su disipación paso a paso.
\end{observacion}

\subsection{Semántica markoviana}

\begin{definicion}[Proceso de Markov discreto sobre $\Delta_n$]\label{def:proceso_markov}
Sea $\mathcal{S} = \{1,\dots,n\}$ un espacio de estados finito. Una \emph{cadena de Markov homogénea} sobre $\mathcal{S}$ con distribución inicial $X_0 \in \Delta_n$ y matriz de transición $K \in \mathcal{S}_n$ es una sucesión de variables aleatorias $(\xi_t)_{t\geq 0}$ sobre $\mathcal{S}$ tal que:
\begin{enumerate}[label=\textnormal{(\roman*)}]
    \item $P(\xi_0 = i) = X_{0,i}$ para todo $i=1,\dots,n$;
    \item $P(\xi_{t+1}=j \mid \xi_t=i,\, \xi_{t-1},\dots,\xi_0) = K_{ij}$ para todo $t\geq 0$ y todo $i,j$.
\end{enumerate}
El \emph{estado informacional en el tiempo $t$} es el vector
\[
X_t = (P(\xi_t=1),\dots,P(\xi_t=n)) \in \Delta_n,
\]
que representa la distribución marginal de $\xi_t$. Por la Proposición~\ref{prop:simplex_invariante}, $X_t\in\Delta_n$ para todo $t\ge 0$. Si además $K_{ij}>0$ para todo $i,j$, entonces $X_t\in\Delta_n^\circ$ para todo $t\ge 1$.
\end{definicion}

\begin{proposicion}[El operador de Markov como corrección local dependiente del estado]\label{prop:markov_como_CI}
Sea $K \in \mathcal{S}_n$ y sea $X \in \Delta_n^\circ$. El vector de corrección local $Y^K(X)$ definido en la Proposición~\ref{prop:markov_es_local} satisface $Y^K(X) \in C(X)$ y
\[
T^{\mathrm{loc}}_{Y^K(X)}(X) = XK.
\]
\end{proposicion}

\begin{proof}
Probado en la Proposición~\ref{prop:markov_es_local}.
\end{proof}

\begin{observacion}[Carácter local de Markov frente al carácter global de Bayes]\label{obs:local_vs_global_markov_bayes}
La corrección $Y^K(X)$ depende del estado base $X$: es una corrección \emph{local} en el sentido de la Sección~B. La verosimilitud bayesiana $Y^\ell$ no depende de $X$: es una corrección \emph{global} en el sentido de la Sección~C. La Proposición~\ref{prop:markov_es_local} demuestra que esta asimetría es estructural: Markov mezcla coordenadas mediante $K$; Bayes escala coordenadas por factores fijos.
\end{observacion}

\begin{definicion}[Distribución estacionaria en CI]\label{def:estacionaria_CI}
Sea $K \in \mathcal{S}_n$. Una distribución $\pi \in \Delta_n^\circ$ se llama \emph{estacionaria} para $K$ si
\[
\pi K = \pi.
\]
Equivalentemente en CI: $\pi$ es estacionaria si y sólo si $Y^K(\pi) = 0$, pues $Y^K_i(\pi) = (\pi K)_i/\pi_i - 1 = 0$ para todo $i$.
\end{definicion}

\begin{proposicion}[Disipación de la divergencia KL hacia la estacionaria]\label{prop:disipacion_KL_markov}
Sea $K \in \mathcal{S}_n$ y sea $\pi \in \Delta_n^\circ$ su distribución estacionaria. Sea $E_\pi(X) = D_{\mathrm{KL}}(\pi \| X)$ el funcional de la Sección~D. Entonces, para todo $X \in \Delta_n^\circ$,
\[
E_\pi(XK) \leq E_\pi(X).
\]
Es decir, un paso del operador de Markov no incrementa la divergencia KL hacia la estacionaria.
\end{proposicion}

\begin{proof}
Como $\pi$ es estacionaria, $\pi_j = \sum_i \pi_i K_{ij}$ para todo $j$. Por tanto,
\[
D_{\mathrm{KL}}(\pi \| XK) = \sum_j \pi_j \log \frac{\pi_j}{(XK)_j} = \sum_j \left(\sum_i \pi_i K_{ij}\right) \log \frac{\sum_i \pi_i K_{ij}}{\sum_i X_i K_{ij}}.
\]
Para cada $j$ fijo, los pesos $w_{ij} := \pi_i K_{ij}/\pi_j$ satisfacen $w_{ij} \geq 0$ y $\sum_i w_{ij}=1$. Aplicando la desigualdad de Jensen a la función cóncava $\log$,
\[
\log \frac{\sum_i \pi_i K_{ij}}{\sum_i X_i K_{ij}} \leq \sum_i w_{ij} \log \frac{\pi_i}{X_i}.
\]
Multiplicando por $\pi_j$ y sumando en $j$,
\[
D_{\mathrm{KL}}(\pi\|XK) \leq \sum_j \sum_i \pi_i K_{ij} \log \frac{\pi_i}{X_i} = \sum_i \pi_i \log\frac{\pi_i}{X_i} \underbrace{\sum_j K_{ij}}_{=1} = D_{\mathrm{KL}}(\pi\|X).
\]
\end{proof}

\begin{proposicion}[Trayectoria markoviana como trayectoria local de CI]\label{prop:trayectoria_markov_CI}
Sea $K \in \mathcal{S}_n$ y sea $X_0 \in \Delta_n$. La trayectoria markoviana $X_{t+1} = X_t K$ satisface $X_t\in\Delta_n$ para todo $t\ge 0$ (Proposición~\ref{prop:simplex_invariante}). Cuando $X_t\in\Delta_n^\circ$, cada paso puede escribirse como
\[
X_{t+1} = T^{\mathrm{loc}}_{Y^K(X_t)}(X_t),
\]
y en forma cerrada,
\[
X_t = X_0 K^t.
\]
Si además $K$ es primitiva (existe $m\geq 1$ con $K^m_{ij}>0$ para todo $i,j$), entonces $K$ admite una única distribución estacionaria $\pi \in \Delta_n^\circ$, y para todo $X_0\in\Delta_n^\circ$,
\[
E_\pi(X_t) = D_{\mathrm{KL}}(\pi\|X_t) \searrow 0 \quad (t\to\infty),
\]
con convergencia garantizada por la Proposición~\ref{prop:disipacion_KL_markov} y la ergodicidad de $K$.
\end{proposicion}

\begin{proof}
La invariancia de $\Delta_n$ es la Proposición~\ref{prop:simplex_invariante}. Para $X_t\in\Delta_n^\circ$, la representación como operador local es la Proposición~\ref{prop:markov_como_CI}. La forma cerrada $X_t = X_0 K^t$ se obtiene por inducción: si $X_t = X_0 K^t$, entonces $X_{t+1} = X_t K = X_0 K^{t+1}$. La convergencia $E_\pi(X_t)\to 0$ se sigue de la Proposición~\ref{prop:disipacion_KL_markov} y del teorema de Perron--Frobenius aplicado a matrices primitivas, que garantiza $X_t \to \pi$.
\end{proof}

\subsection{Cadenas absorbentes y frontera del símplex}

Las cadenas absorbentes extienden el marco markoviano hacia la frontera $\partial\Delta_n$. Los estados absorbentes son vértices canónicos $\delta_a\in\partial\Delta_n$ que son puntos fijos de $L_K$; los estados transitorios viven inicialmente en el interior pero convergen a la frontera. Esta subsección formaliza esa extensión dentro del lenguaje de CI, fundándose en la Proposición~\ref{prop:simplex_invariante} y en la Proposición~\ref{prop:trayectoria_markov_CI}.

\begin{definicion}[Estado absorbente en CI]\label{def:estado_absorbente}
Sea $K\in\mathcal{S}_n$. Un estado $a\in\{1,\dots,n\}$ se llama \emph{absorbente} si $K_{aa}=1$ (y por tanto $K_{aj}=0$ para $j\neq a$). El vértice canónico $\delta_a\in\partial\Delta_n$ satisface
\[
\delta_a K = \delta_a,
\]
es decir, $\delta_a$ es un punto fijo de $L_K$ en la frontera del símplex.

En una cadena absorbente se distinguen los estados absorbentes $\mathcal{A}$ y los estados transitorios $\mathcal{T}$, entendidos como aquellos desde los cuales la cadena alcanza algún estado absorbente con probabilidad uno. La cadena se dice \emph{absorbente} si todo estado no absorbente pertenece a $\mathcal{T}$, es decir, si desde cualquier estado no absorbente existe una trayectoria con probabilidad positiva hacia algún estado de $\mathcal{A}$.
\end{definicion}

\begin{observacion}[Estados absorbentes y frontera de $\Delta_n$]\label{obs:absorbentes_frontera}
Los estados absorbentes viven en la frontera $\partial\Delta_n = \Delta_n\setminus\Delta_n^\circ$. Por eso la corrección local $Y^K_i(X) = (XK)_i/X_i - 1$ no está definida en $\delta_a$: la coordenada $a$ vale $1$ y las demás valen $0$, haciendo los cocientes indefinidos en las coordenadas nulas. La trayectoria $X_t = X_0 K^t$ sí está bien definida para todo $t$ por la Proposición~\ref{prop:simplex_invariante}, incluso cuando $X_t$ alcanza la frontera. La representación como corrección local (Proposición~\ref{prop:markov_como_CI}) aplica solo mientras $X_t\in\Delta_n^\circ$.
\end{observacion}

\begin{definicion}[Forma canónica de una cadena absorbente]\label{def:forma_canonica_absorbente}
Sea $K\in\mathcal{S}_n$ una cadena absorbente con $r$ estados absorbentes y $s=n-r$ estados transitorios. Reordenando los estados con los transitorios primero y los absorbentes al final, la matriz $K$ toma la \emph{forma canónica}
\[
K =
\begin{pmatrix}
Q & R \\
0 & I_r
\end{pmatrix},
\]
donde $Q\in\mathbb{R}^{s\times s}$ es la submatriz de transición entre estados transitorios, $R\in\mathbb{R}^{s\times r}$ es la submatriz de transición de estados transitorios a absorbentes, $0$ es la matriz nula $r\times s$, e $I_r$ es la identidad $r\times r$.
\end{definicion}

\begin{proposicion}[Matriz fundamental y absorción]\label{prop:matriz_fundamental}
Sea $K$ una cadena absorbente en forma canónica (Definición~\ref{def:forma_canonica_absorbente}). Entonces:
\begin{enumerate}[label=\textnormal{(\roman*)}]
    \item La serie $\sum_{t=0}^\infty Q^t$ converge y su suma es la \emph{matriz fundamental}
    \[
    N := (I_s - Q)^{-1}.
    \]
    \item La entrada $N_{ij}$ es el número esperado de veces que la cadena visita el estado transitorio $j$ antes de ser absorbida, dado que parte del estado transitorio $i$.
    \item La \emph{matriz de probabilidades de absorción} es
    \[
    B := NR \in \mathbb{R}^{s\times r},
    \]
    donde $B_{ia}$ es la probabilidad de ser absorbido en el estado absorbente $a$, dado que la cadena parte del estado transitorio $i$.
    \item El \emph{tiempo esperado de absorción} desde el estado transitorio $i$ es
    \[
    \tau_i = (N\mathbf{1})_i = \sum_{j=1}^s N_{ij},
    \]
    donde $\mathbf{1}$ es el vector de unos.
\end{enumerate}
\end{proposicion}

\begin{proof}
Como $K$ es absorbente, todos los valores propios de $Q$ tienen módulo estrictamente menor que $1$. Por tanto $(I_s-Q)$ es invertible y $Q^t\to 0$ cuando $t\to\infty$, lo que garantiza la convergencia de $\sum_{t=0}^\infty Q^t = (I_s-Q)^{-1} = N$.

Para (ii): la entrada $(Q^t)_{ij}$ es la probabilidad de estar en el estado transitorio $j$ en el tiempo $t$ habiendo partido de $i$ sin haber sido absorbido. Sumando en $t$, $N_{ij} = \sum_{t=0}^\infty (Q^t)_{ij}$ es el número esperado de visitas a $j$ antes de absorción.

Para (iii): $R_{ia}$ es la probabilidad de pasar del estado transitorio $i$ al absorbente $a$ en un paso. Condicionando sobre el primer paso, la probabilidad total de absorción en $a$ desde $i$ es $(NR)_{ia}$. Como la cadena es absorbente, $B\mathbf{1} = \mathbf{1}$ (la cadena es absorbida con probabilidad $1$).

Para (iv): el número esperado de pasos hasta la absorción desde $i$ es la suma de las visitas esperadas a todos los estados transitorios, que es $\sum_j N_{ij} = (N\mathbf{1})_i$.
\end{proof}

\begin{observacion}[Lectura en CI: absorción como convergencia a la frontera]\label{obs:absorcion_CI}
En el lenguaje de CI, la Proposición~\ref{prop:matriz_fundamental} describe cómo la trayectoria $X_t = X_0 K^t$ converge a un punto de la frontera $\partial\Delta_n$. Si $X_0 = \delta_i$ (estado transitorio $i$), entonces
\[
\lim_{t\to\infty} X_t = \sum_{a} B_{ia}\,\delta_a,
\]
que es una combinación convexa de los vértices absorbentes con pesos dados por las probabilidades de absorción $B_{ia}$. Este límite pertenece a la cara del símplex generada por los vértices absorbentes.
\end{observacion}

\begin{definicion}[Ecuaciones de primer paso]\label{def:primer_paso}
Sea $K$ una cadena absorbente con estados absorbentes $\mathcal{A}$ y estados transitorios $\mathcal{T}$. Sea $a\in\mathcal{A}$ un estado absorbente fijo. La \emph{probabilidad de absorción en $a$ desde $i$} es $h_i^{(a)} := P_i(T_a < T_{\mathcal{A}\setminus\{a\}})$, donde $T_b$ denota el primer tiempo de llegada al estado $b$.

Estas probabilidades satisfacen las \emph{ecuaciones de primer paso}:
\[
h_i^{(a)} = \sum_{j\in\mathcal{T}} K_{ij}\,h_j^{(a)} + K_{ia}, \qquad i\in\mathcal{T},
\]
con condiciones de frontera $h_a^{(a)} = 1$ y $h_{a'}^{(a)} = 0$ para $a'\in\mathcal{A}$, $a'\neq a$.

En forma matricial: $h^{(a)} = Qh^{(a)} + R_{\cdot a}$, es decir,
\[
h^{(a)} = (I_s-Q)^{-1}R_{\cdot a} = NR_{\cdot a} = B_{\cdot a},
\]
que es la columna $a$ de la matriz $B = NR$.
\end{definicion}

\begin{corolario}[Ruina del jugador en CI]\label{cor:ruina_jugador}
Sea $\mathcal{S} = \{0,1,\dots,N\}$ con estados absorbentes $\{0,N\}$ y estados transitorios $\{1,\dots,N-1\}$. En cada paso, desde el estado $i\in\{1,\dots,N-1\}$, la cadena transita a $i+1$ con probabilidad $p$ y a $i-1$ con probabilidad $q=1-p$.

La probabilidad de absorción en $N$ (no arruinarse) desde el estado $i$ satisface las ecuaciones de primer paso
\[
h_i = ph_{i+1} + qh_{i-1}, \qquad 1\le i\le N-1,
\]
con $h_0=0$ y $h_N=1$. La solución es:
\[
h_i =
\begin{cases}
\dfrac{1-(q/p)^i}{1-(q/p)^N} & \text{si } p\neq q, \\[1em]
\dfrac{i}{N} & \text{si } p=q=\tfrac{1}{2}.
\end{cases}
\]

El tiempo esperado de absorción desde $i$ satisface
\[
\tau_i = 1 + p\tau_{i+1} + q\tau_{i-1}, \qquad 1\le i\le N-1,
\]
con $\tau_0=\tau_N=0$. La solución es:
\[
\tau_i =
\begin{cases}
\dfrac{1}{q-p}\!\left[i - N\dfrac{1-(q/p)^i}{1-(q/p)^N}\right] & \text{si } p\neq q, \\[1em]
i(N-i) & \text{si } p=q=\tfrac{1}{2}.
\end{cases}
\]

En CI: la trayectoria informacional converge a la cara generada por $\{\delta_0,\delta_N\}$ de $\partial\Delta_{N+1}$,
\[
\lim_{t\to\infty}\delta_i K^t = (1-h_i)\delta_0 + h_i\delta_N,
\]
con pesos $(1-h_i,\,h_i)$.
\end{corolario}

\begin{proof}
Las ecuaciones de primer paso para $h_i$ son la Definición~\ref{def:primer_paso} aplicada a la cadena de ruina del jugador, que tiene $Q_{i,i+1}=p$ y $Q_{i,i-1}=q$ para $1\le i\le N-1$ (ceros en el resto). La solución de la recurrencia lineal de segundo orden $h_i = ph_{i+1}+qh_{i-1}$, equivalente a $p(h_{i+1}-h_i) = q(h_i-h_{i-1})$, es geométrica si $p\neq q$ y lineal si $p=q$, con constantes determinadas por las condiciones de frontera $h_0=0$, $h_N=1$.

Las ecuaciones para $\tau_i$ son la Definición~\ref{def:primer_paso} aplicada al tiempo esperado, que satisface $\tau_i = 1 + p\tau_{i+1}+q\tau_{i-1}$ con $\tau_0=\tau_N=0$. La solución se obtiene por variación de parámetros sobre la solución homogénea.

La lectura en CI se sigue de la Observación~\ref{obs:absorcion_CI}: $\lim_{t\to\infty}\delta_i K^t = B_{i,0}\delta_0 + B_{i,N}\delta_N = (1-h_i)\delta_0 + h_i\delta_N$.
\end{proof}

\subsection{Clasificación de estados en CI}

La clasificación de estados de una cadena de Markov ---comunicación, recurrencia, transitividad y periodicidad--- se puede formular completamente dentro de CI usando únicamente el operador $L_K$ y la trayectoria $X_t = X_0 K^t$ ya establecidos. El objeto clave es la potencia $K^t$, cuyas entradas $(K^t)_{ij}$ son accesibles como coordenadas de la trayectoria $\delta_i K^t \in \Delta_n$.

\begin{definicion}[Accesibilidad y comunicación]\label{def:accesibilidad}
Sea $K\in\mathcal{S}_n$. Se dice que el estado $j$ es \emph{accesible} desde el estado $i$, y se escribe $i\to j$, si existe $t\ge 0$ tal que
\[
(K^t)_{ij} > 0.
\]
En el lenguaje de CI, esto equivale a que la trayectoria $\delta_i K^t$ tenga masa positiva en la coordenada $j$ para algún $t$, es decir,
\[
(\delta_i K^t)_j > 0 \quad \text{para algún } t\ge 0.
\]

Dos estados $i$ y $j$ se dicen \emph{comunicantes}, y se escribe $i\leftrightarrow j$, si $i\to j$ y $j\to i$.
\end{definicion}

\begin{proposicion}[Comunicación es relación de equivalencia]\label{prop:comunicacion_equivalencia}
La relación $\leftrightarrow$ es una relación de equivalencia sobre $\{1,\dots,n\}$. Sus clases de equivalencia se llaman \emph{clases comunicantes}.
\end{proposicion}

\begin{proof}
\emph{Reflexividad}: $(K^0)_{ii} = 1 > 0$, luego $i\to i$ para todo $i$.

\emph{Simetría}: por definición de $\leftrightarrow$.

\emph{Transitividad}: si $i\to j$ y $j\to k$, existen $s,t\ge 0$ con $(K^s)_{ij}>0$ y $(K^t)_{jk}>0$. Entonces
\[
(K^{s+t})_{ik} = \sum_m (K^s)_{im}(K^t)_{mk} \ge (K^s)_{ij}(K^t)_{jk} > 0,
\]
luego $i\to k$. Si además $j\to i$ y $k\to j$, entonces $k\to i$ por el mismo argumento.
\end{proof}

\begin{definicion}[Clase cerrada]\label{def:clase_cerrada}
Una clase comunicante $C\subset\{1,\dots,n\}$ se dice \emph{cerrada} si
\[
K_{ij} = 0 \quad \text{para todo } i\in C,\ j\notin C.
\]
En CI: $C$ es cerrada si y sólo si $L_K(\delta_i) = \delta_i K$ tiene soporte contenido en $C$ para todo $i\in C$, es decir, la trayectoria desde cualquier estado de $C$ permanece en $C$.
\end{definicion}

\begin{observacion}[Clases cerradas y caras del símplex]\label{obs:clases_caras}
Una clase cerrada $C$ corresponde a una \emph{cara} del símplex $\Delta_n$: el subconjunto
\[
F_C := \{X\in\Delta_n : X_i=0 \ \forall i\notin C\}
\]
es cerrado bajo $L_K$, pues si $X\in F_C$ entonces $(XK)_j = \sum_i X_i K_{ij} = 0$ para todo $j\notin C$ (ya que $K_{ij}=0$ cuando $i\in C$ y $j\notin C$). Los estados absorbentes de la subsección anterior son el caso particular $|C|=1$.
\end{observacion}

\begin{definicion}[Estados recurrentes y transitorios]\label{def:recurrente_transitorio}
Sea $i\in\{1,\dots,n\}$. El estado $i$ se llama \emph{recurrente} si
\[
\sum_{t=1}^\infty (K^t)_{ii} = +\infty,
\]
y \emph{transitorio} si
\[
\sum_{t=1}^\infty (K^t)_{ii} < +\infty.
\]

En CI: $(K^t)_{ii} = (\delta_i K^t)_i$ es la masa que la trayectoria $\delta_i K^t$ devuelve a la coordenada $i$ en el tiempo $t$. El estado $i$ es recurrente si la masa acumulada de retorno es infinita, y transitorio si es finita.
\end{definicion}

\begin{proposicion}[Recurrencia y transitividad son propiedades de clase]\label{prop:recurrencia_clase}
Si $i\leftrightarrow j$, entonces $i$ es recurrente si y sólo si $j$ es recurrente. En particular, todos los estados de una clase comunicante son simultáneamente recurrentes o simultáneamente transitorios.
\end{proposicion}

\begin{proof}
Si $i\leftrightarrow j$, existen $s,t\ge 0$ con $\alpha:=(K^s)_{ij}>0$ y $\beta:=(K^t)_{ji}>0$. Para todo $m\ge 0$,
\[
(K^{s+m+t})_{ii} \ge (K^s)_{ij}(K^m)_{jj}(K^t)_{ji} = \alpha\beta (K^m)_{jj}.
\]
Sumando en $m$: $\sum_m (K^{s+m+t})_{ii} \ge \alpha\beta \sum_m (K^m)_{jj}$. Si $j$ es recurrente, el lado derecho diverge, luego $i$ es recurrente. Por simetría ($i\leftrightarrow j$ implica $j\leftrightarrow i$), también se tiene la implicación inversa.
\end{proof}

\begin{proposicion}[Caracterización de clases recurrentes]\label{prop:clases_recurrentes}
Una clase comunicante $C$ es recurrente si y sólo si es cerrada. En particular:
\begin{enumerate}[label=\textnormal{(\roman*)}]
    \item Todo estado absorbente (clase cerrada de tamaño uno) es recurrente.
    \item Todo estado en una clase no cerrada es transitorio.
    \item En una cadena con espacio de estados finito, existe al menos una clase recurrente.
\end{enumerate}
\end{proposicion}

\begin{proof}
Si $C$ es cerrada, la trayectoria desde $i\in C$ permanece en $C$ (Definición~\ref{def:clase_cerrada}). Como $C$ es finita, la masa $(\delta_i K^t)_i$ no puede tender a cero sin retornar, y $\sum_t (K^t)_{ii} = +\infty$.

Si $C$ no es cerrada, existe $i\in C$ y $j\notin C$ con $K_{ij}>0$. Una vez que la trayectoria sale de $C$ no puede regresar (pues $j\notin C$ y $C$ no es accesible desde fuera por definición de clase comunicante). Luego $\sum_t (K^t)_{ii}<+\infty$ y $i$ es transitorio.

Para (iii): en un espacio finito, la masa de la trayectoria debe acumularse en algún lugar; ese lugar es necesariamente una clase cerrada, que es recurrente por lo anterior.
\end{proof}

\begin{definicion}[Periodo de un estado]\label{def:periodo}
Sea $i\in\{1,\dots,n\}$. El \emph{periodo} del estado $i$ es
\[
d(i) := \gcd\{t\ge 1 : (K^t)_{ii} > 0\}.
\]
Si el conjunto es vacío, se conviene $d(i)=+\infty$. En CI: $d(i)$ es el máximo común divisor de los tiempos $t$ para los cuales la trayectoria $\delta_i K^t$ devuelve masa positiva a la coordenada $i$.

El estado $i$ se llama \emph{aperiódico} si $d(i)=1$.
\end{definicion}

\begin{proposicion}[El periodo es invariante en la clase]\label{prop:periodo_clase}
Si $i\leftrightarrow j$, entonces $d(i)=d(j)$.
\end{proposicion}

\begin{proof}
Sean $s,t\ge 0$ con $(K^s)_{ij}>0$ y $(K^t)_{ji}>0$. Si $(K^m)_{jj}>0$, entonces
\[
(K^{s+m+t})_{ii} \ge (K^s)_{ij}(K^m)_{jj}(K^t)_{ji} > 0,
\]
luego $d(i)\mid s+m+t$. Como también $(K^{s+t})_{ii}\ge (K^s)_{ij}(K^t)_{ji}>0$, se tiene $d(i)\mid s+t$. Por tanto $d(i)\mid m$ para todo $m$ con $(K^m)_{jj}>0$, lo que implica $d(i)\mid d(j)$. Por simetría $d(j)\mid d(i)$, luego $d(i)=d(j)$.
\end{proof}

\begin{definicion}[Cadena irreducible y ergódica]\label{def:irreducible_ergodica}
La cadena $K$ se llama \emph{irreducible} si tiene una única clase comunicante, es decir, si $i\leftrightarrow j$ para todo par $i,j\in\{1,\dots,n\}$.

La cadena es \emph{ergódica} si es irreducible y aperiódica ($d(i)=1$ para todo $i$). Para cadenas ergódicas, la Proposición~\ref{prop:trayectoria_markov_CI} garantiza que $X_t\to\pi$ para todo $X_0\in\Delta_n$, donde $\pi$ es la única distribución estacionaria.
\end{definicion}

\begin{proposicion}[Convergencia para cadenas ergódicas desde CI]\label{prop:convergencia_ergodica}
Sea $K\in\mathcal{S}_n$ ergódica con distribución estacionaria $\pi\in\Delta_n^\circ$. Entonces para todo $X_0\in\Delta_n^\circ$,
\[
E_\pi(X_t) = D_{\mathrm{KL}}(\pi\|X_t) \searrow 0 \quad (t\to\infty).
\]
En particular, $X_t \to \pi$.
\end{proposicion}

\begin{proof}
Por la Proposición~\ref{prop:disipacion_KL_markov}, la sucesión $E_\pi(X_t)$ es decreciente y acotada por abajo por $0$ (Proposición~\ref{prop:minimoGlobalKL}), luego converge a algún $L\ge 0$. Supóngase que $L>0$; entonces $X_t$ no converge a $\pi$. Como $\Delta_n$ es compacto, existe una subsucesión $X_{t_k}\to X^*\neq\pi$. Por continuidad de $L_K$ y la disipación, $E_\pi(X^*K)\le E_\pi(X^*)=L$. Pero la ergodicidad de $K$ implica que $(K^t)_{ij}>0$ para $t$ suficientemente grande, de modo que la trayectoria desde $X^*$ se acerca a $\pi$, contradiciendo $E_\pi(X^*)=L>0$. Por tanto $L=0$ y $X_t\to\pi$.
\end{proof}

\begin{observacion}[Alcance de la convergencia]\label{obs:alcance_convergencia}
La Proposición~\ref{prop:convergencia_ergodica} garantiza convergencia pero no cuantifica su velocidad. La tasa de convergencia depende de la \emph{brecha espectral} de $K$, definida como $1-|\lambda_2|$ donde $\lambda_2$ es el segundo valor propio de $K$ en módulo. Esa cuantificación requiere álgebra lineal espectral que queda fuera del alcance de este documento.
\end{observacion}

\begin{observacion}[No conmutatividad entre Bayes y Markov]\label{obs:no_conmutacion}
En general, el operador bayesiano $T_{Y^\ell}$ y el operador de Markov $L_K$ \emph{no conmutan}. Las dos composiciones posibles son
\[
T_{Y^\ell}(XK)_j = \frac{(XK)_j \ell_j}{\sum_m (XK)_m \ell_m}
\qquad\text{y}\qquad
(T_{Y^\ell}(X)\,K)_j = \frac{\sum_i X_i\ell_i K_{ij}}{\sum_k X_k\ell_k},
\]
que son expresiones distintas porque la primera normaliza \emph{después} de aplicar $K$, mientras que la segunda aplica $K$ a un vector ya normalizado por $\ell$. La diferencia entre ambos estados mide el efecto de intercambiar el orden de actualización informacional y transición estocástica. Determinar las condiciones exactas sobre $(K,\ell)$ bajo las cuales estas dos expresiones coinciden para todo $X\in\Delta_n^\circ$ queda como problema abierto dentro del marco de CI.
\end{observacion}

\begin{observacion}[Tabla de correspondencias CI $\leftrightarrow$ Bayes/Markov]\label{obs:tabla_correspondencias}
El siguiente cuadro resume cómo los objetos del marco de CI se corresponden con los de los paradigmas clásicos.

\medskip
\begin{center}
\begin{tabular}{lll}
\toprule
\textbf{Control Informacional} & \textbf{Bayes clásico} & \textbf{Markov clásico} \\
\midrule
$X \in \Delta_n^\circ$ & distribución a priori & --- \\
$X \in \Delta_n$ & --- & distribución en el tiempo $t$ \\
$Y \in C_{\mathrm{glob}}$, $1+Y_i = \ell_i$ & verosimilitud $P(E_j\mid H_i)$ & --- \\
$T_{Y^\ell}(X)$ (Sección~C) & posterior $P(H_i\mid E_j)$ & --- \\
$Y^K(X) \in C(X)$ & --- & corrección local de un paso \\
$T^{\mathrm{loc}}_{Y^K(X)}(X)$ (Sección~B) & --- & $XK$, un paso de Markov \\
$\pi$ : $Y^K(\pi)=0$ & --- & distribución estacionaria \\
$E_\pi(X) = D_{\mathrm{KL}}(\pi\|X)$ (Sección~D) & distancia al posterior & energía hacia la estacionaria \\
$\mathcal{G}$, composición & actualizaciones secuenciales & no representa Markov globalmente \\
\bottomrule
\end{tabular}
\end{center}
\end{observacion}

\begin{ejemplo}[Diagnóstico médico bayesiano en CI]\label{ej:diagnostico}
Sea $n=2$, $H_1 = \{\text{enfermo}\}$, $H_2 = \{\text{sano}\}$. Prior: $X=(0.01, 0.99)$. Dato $e$: prueba positiva, con $\ell_1 = P(e\mid H_1)=0.99$ y $\ell_2=P(e\mid H_2)=0.05$. La corrección global es $1+Y_i^\ell = \ell_i$, y el posterior es
\[
T_{Y^\ell}(X)_1 = \frac{0.01\times 0.99}{0.01\times 0.99 + 0.99\times 0.05} = \frac{0.0099}{0.0099+0.0495} \approx 0.167.
\]
\end{ejemplo}

\subsection{Modelos continuos como estados informacionales discretos}

El marco de las Secciones A--D opera exclusivamente sobre $\Delta_n$ con $n$ finito. Sin embargo, ciertos modelos probabilísticos clásicos ---como los modelos conjugados Beta-Binomial y Normal-Normal--- tienen parámetros en espacios continuos. Esta subsección muestra cómo ingresar dichos modelos al marco de CI sin salir de él: el espacio continuo se convierte en un espacio de probabilidad $(\Omega,\mathcal{F},P_0)$ concreto (Definición~\ref{def:espacio_prob}), la cuadrícula define una partición medible finita no degenerada (Definición~\ref{def:particion}), y las probabilidades de los eventos de esa partición producen el estado informacional en $\Delta_n^\circ$ (Definición~\ref{def:estado} y Proposición~\ref{prop:estado_simplex}). La verosimilitud entra como corrección global (Definición~\ref{def:Cglob}) y la actualización es el operador $T_{Y^\ell}$ de la Sección~C (Proposición~\ref{prop:bayes_como_TY}).

\begin{definicion}[Partición paramétrica uniforme]\label{def:particion_parametrica}
Sea $\Omega = (a,b) \subset \mathbb{R}$ un intervalo abierto y sea $\mathcal{F}$ la $\sigma$-álgebra de Borel en $\Omega$. Sea $P_0$ una medida de probabilidad absolutamente continua sobre $(\Omega,\mathcal{F})$ con densidad $f:\Omega\to(0,\infty)$.

Dado $n\geq 1$, sea $h:=(b-a)/n$ el paso de discretización y defínase la \emph{partición paramétrica uniforme} de paso $h$ como la familia
\[
\Pi_n := \{E_1,\dots,E_n\},
\qquad
E_i := \left(a+(i-1)h,\; a+ih\right], \quad i=1,\dots,n.
\]
Esta es una partición medible finita de $\Omega$ en el sentido de la Definición~\ref{def:particion}. Como $f>0$ en $\Omega$, se tiene $P_0(E_i)>0$ para todo $i$, de modo que $\Pi_n$ es no degenerada.

El \emph{punto representativo} del intervalo $E_i$ es el punto medio
\[
\theta_i := a + \left(i-\tfrac{1}{2}\right)h.
\]
El \emph{estado informacional previo} asociado a $(\Omega,\mathcal{F},P_0,\Pi_n)$ es, por la Proposición~\ref{prop:estado_simplex},
\[
X_{0,i} := P_0(E_i) = \int_{E_i} f(\theta)\,d\theta, \qquad i=1,\dots,n,
\]
y $X_0 \in \Delta_n^\circ$.
\end{definicion}

\begin{observacion}[Relación entre $X_{0,i}$ y $f(\theta_i)$]\label{obs:aprox_midpoint}
La definición exacta es $X_{0,i} = P_0(E_i) = \int_{E_i} f(\theta)\,d\theta$. Si $f\in C^2(\Omega)$, la regla del punto medio da
\[
X_{0,i} = f(\theta_i)\,h + O(h^3),
\]
donde el error local es de orden $O(h^3)$ por intervalo. Esto motiva introducir el \emph{prior discreto por punto medio}
\[
\widetilde X_{0,i} := \frac{f(\theta_i)}{\displaystyle\sum_{k=1}^n f(\theta_k)},
\qquad i=1,\dots,n,
\]
que pertenece a $\Delta_n^\circ$ y aproxima $X_{0,i}$ con error $O(h^2)$ en cada coordenada normalizada. En lo que sigue se distinguirá explícitamente cuándo se usa la masa exacta $X_{0,i}=P_0(E_i)$ y cuándo la aproximación $\widetilde X_{0,i}\propto f(\theta_i)$.
\end{observacion}

\begin{definicion}[Verosimilitud como corrección global admisible]\label{def:verosimilitud_global}
Sea $e$ un dato observado y sea $\ell:(a,b)\to(0,\infty)$ una función de verosimilitud, es decir, $\ell(\theta) = P(e\mid \theta)$ para cada valor del parámetro $\theta$. Sobre la partición $\Pi_n$, la verosimilitud del dato $e$ bajo el evento $E_i$ se define como
\[
\ell_i := \ell(\theta_i),
\qquad i=1,\dots,n.
\]
Como $\ell(\theta)>0$ para todo $\theta\in(a,b)$, se tiene $\ell_i>0$ para todo $i$. Definiendo $Y^\ell \in \mathbb{R}^n$ por $1+Y^\ell_i := \ell_i$, se obtiene $Y^\ell \in C_{\mathrm{glob}}$ (Definición~\ref{def:Cglob}).
\end{definicion}

\begin{observacion}[Fundamento en la Sección B: cambio de medida normalizado]\label{obs:verosimilitud_cambio_medida}
La corrección $Y^\ell$ induce un cambio de medida sobre $(\Omega,\mathcal{F},P_0)$ en el sentido de la Proposición~\ref{prop:QYprobabilidad}. La densidad de Radon-Nikodym de la medida posterior $Q_{Y^\ell}$ respecto de $P_0$ no es $\ell(\theta_i)$ directamente, sino su versión normalizada. Definiendo la constante de normalización
\[
Z_n := \sum_{k=1}^n P_0(E_k)\,\ell(\theta_k),
\]
la función
\[
\widetilde h(\omega) := \frac{\ell(\theta_i)}{Z_n}, \qquad \omega\in E_i,
\]
satisface $\int_\Omega \widetilde h\,dP_0 = 1$ y define una medida de probabilidad $Q_{Y^\ell}(A) = \int_A \widetilde h\,dP_0$. El estado informacional discreto de $Q_{Y^\ell}$ sobre $\Pi_n$ es entonces
\[
Q_{Y^\ell}(E_i) = \frac{P_0(E_i)\,\ell(\theta_i)}{Z_n} = T_{Y^\ell}(X_0)_i,
\]
donde la última igualdad se sigue de la Proposición~\ref{prop:compatibilidadCambioMedida}, y la normalización $Z_n$ queda absorbida por el denominador de $T_{Y^\ell}$.
\end{observacion}

\begin{proposicion}[Actualización bayesiana paramétrica en CI]\label{prop:bayes_parametrico}
Sea $(\Omega,\mathcal{F},P_0,\Pi_n)$ una partición paramétrica uniforme con estado previo $X_0\in\Delta_n^\circ$ y sea $Y^\ell\in C_{\mathrm{glob}}$ la corrección de verosimilitud de la Definición~\ref{def:verosimilitud_global}. Entonces
\[
T_{Y^\ell}(X_0)_i
= \frac{P_0(E_i)\,\ell(\theta_i)}{\displaystyle\sum_{k=1}^n P_0(E_k)\,\ell(\theta_k)},
\qquad i=1,\dots,n,
\]
y $T_{Y^\ell}(X_0)\in\Delta_n^\circ$. Este vector es la distribución posterior discreta de $\theta$ dado el dato $e$, obtenida aplicando la regla de Bayes discreta sobre la partición $\Pi_n$.

Además, si $f$ y $\ell$ son continuas en $\Omega$, cuando $n\to\infty$ (con $h=(b-a)/n\to 0$),
\[
\sum_{k=1}^n P_0(E_k)\,\ell(\theta_k)
\;\longrightarrow\;
\int_a^b f(\theta)\,\ell(\theta)\,d\theta,
\]
y la distribución discreta $T_{Y^\ell}(X_0)$ converge a la distribución posterior continua con densidad proporcional a $f(\theta)\ell(\theta)$.
\end{proposicion}

\begin{proof}
La primera parte es aplicación directa de la Proposición~\ref{prop:bayes_como_TY}: dado que $X_{0,i}=P_0(E_i)>0$ y $1+Y^\ell_i=\ell_i>0$, el operador global da
\[
T_{Y^\ell}(X_0)_i = \frac{X_{0,i}\,\ell_i}{\sum_{k=1}^n X_{0,k}\,\ell_k} = \frac{P_0(E_i)\,\ell(\theta_i)}{\sum_{k=1}^n P_0(E_k)\,\ell(\theta_k)}.
\]
La pertenencia a $\Delta_n^\circ$ se sigue de la Proposición~\ref{prop:interiorGlobal}.

Para la convergencia: $P_0(E_k)=\int_{E_k}f(\theta)\,d\theta$ y $\ell(\theta_k)$ es el valor de $\ell$ en el punto medio de $E_k$. Por continuidad de $f\cdot\ell$, la suma $\sum_k P_0(E_k)\,\ell(\theta_k)$ es una suma de Riemann que converge a $\int_a^b f(\theta)\ell(\theta)\,d\theta$ cuando $h\to 0$.
\end{proof}

\begin{proposicion}[Modelo conjugado exponencial como caso particular]\label{prop:conjugado_exponencial_CI}
Sea $f(\theta)\propto\exp(\eta_0\cdot T(\theta))$ un prior de la familia exponencial con parámetro natural $\eta_0\in\mathbb{R}^k$ y estadístico suficiente $T:\Omega\to\mathbb{R}^k$. Sea $\ell(\theta)\propto\exp(\eta_\ell\cdot T(\theta))$ la verosimilitud de $m$ observaciones, donde $\eta_\ell$ es la contribución del dato al parámetro natural.

Usando el prior discreto por punto medio $\widetilde X_{0,i}\propto f(\theta_i)$ (Observación~\ref{obs:aprox_midpoint}), el operador $T_{Y^\ell}$ produce un posterior discreto \emph{exactamente} igual a
\[
T_{Y^\ell}(\widetilde X_0)_i
\;\propto\;
\exp\!\bigl((\eta_0+\eta_\ell)\cdot T(\theta_i)\bigr),
\]
que es la evaluación de la densidad de la familia exponencial con parámetro natural actualizado $\eta_n := \eta_0 + \eta_\ell$ en los puntos medios $\theta_i$.

Si en cambio se usan las masas exactas $X_{0,i}=P_0(E_i)$ (Proposición~\ref{prop:bayes_parametrico}), la misma fórmula vale como aproximación con error $O(h^2)$ en las coordenadas normalizadas, bajo la hipótesis $f\in C^2(\Omega)$.

En ambos casos, cuando $n\to\infty$, el posterior discreto converge al posterior continuo con parámetro $\eta_n$.

Los dos casos conjugados clásicos son instancias directas:
\begin{enumerate}[label=\textnormal{(\roman*)}]
    \item \emph{Beta-Binomial.} Prior $\mathrm{Beta}(\alpha,\beta)$: $T(\theta)=(\log\theta,\log(1-\theta))$, $\eta_0=(\alpha-1,\beta-1)$. Verosimilitud binomial para $x$ éxitos en $m$ ensayos: $\eta_\ell=(x,m-x)$. Posterior:
    \[
    T_{Y^\ell}(\widetilde X_0)_i \;\propto\; \theta_i^{\alpha+x-1}(1-\theta_i)^{\beta+m-x-1},
    \]
    evaluación de la densidad $\mathrm{Beta}(\alpha+x,\,\beta+m-x)$ en $\theta_i$.

    \item \emph{Normal-Normal} (varianza conocida $\sigma^2$). Prior $\mathcal{N}(\mu_0,\tau_0^2)$: $T(\theta)=(\theta,-\theta^2/2)$, $\eta_0=(\mu_0/\tau_0^2,1/\tau_0^2)$. Verosimilitud normal para $m$ observaciones con media muestral $\bar x$: $\eta_\ell=(m\bar x/\sigma^2,m/\sigma^2)$. Posterior:
    \[
    T_{Y^\ell}(\widetilde X_0)_i \;\propto\; \exp\!\left(-\frac{(\theta_i-\mu_n)^2}{2\tau_n^2}\right),
    \]
    evaluación de la densidad $\mathcal{N}(\mu_n,\tau_n^2)$ en $\theta_i$, con
    \[
    \frac{1}{\tau_n^2}=\frac{1}{\tau_0^2}+\frac{m}{\sigma^2},
    \qquad
    \mu_n=\tau_n^2\!\left(\frac{\mu_0}{\tau_0^2}+\frac{m\bar x}{\sigma^2}\right).
    \]
\end{enumerate}
\end{proposicion}

\begin{proof}
Con el prior por punto medio, $\widetilde X_{0,i}\propto f(\theta_i)\propto\exp(\eta_0\cdot T(\theta_i))$ exactamente, sin error de aproximación. La verosimilitud evaluada es $\ell_i=\ell(\theta_i)\propto\exp(\eta_\ell\cdot T(\theta_i))$. Aplicando $T_{Y^\ell}$ y cancelando constantes independientes de $i$ por invariancia de escala (Proposición~\ref{prop:composicionGlobal}):
\[
T_{Y^\ell}(\widetilde X_0)_i
\propto \exp(\eta_0\cdot T(\theta_i))\cdot\exp(\eta_\ell\cdot T(\theta_i))
= \exp\bigl((\eta_0+\eta_\ell)\cdot T(\theta_i)\bigr).
\]
Con las masas exactas $X_{0,i}=P_0(E_i)=f(\theta_i)h+O(h^3)$, el factor $h$ es constante en $i$ y se cancela; el término $O(h^3)$ introduce un error $O(h^2)$ en las coordenadas normalizadas. Los casos (i) y (ii) se obtienen sustituyendo los estadísticos suficientes y parámetros naturales correspondientes; el factor $\binom{m}{x}$ en el caso binomial es constante en $i$ y se cancela por invariancia de escala. La convergencia cuando $n\to\infty$ se sigue de la Proposición~\ref{prop:bayes_parametrico}.
\end{proof}

\begin{observacion}[Universalidad de la actualización en CI]\label{obs:universalidad}
La Proposición~\ref{prop:conjugado_exponencial_CI} muestra que la regla de actualización conjugada ---la suma de parámetros naturales $\eta_n=\eta_0+\eta_\ell$--- es consecuencia directa de la ley de composición multiplicativa del operador $T_{Y^\ell}$ (Proposición~\ref{prop:composicionGlobal}): el producto de exponenciales en el numerador de $T_{Y^\ell}(X_0)$ corresponde exactamente a la suma de parámetros naturales. La estructura de CI no impone un modelo particular; cualquier par prior-verosimilitud de la familia exponencial produce un posterior en la misma familia, y su discretización sobre $\Pi_n$ es un estado admisible en $\Delta_n^\circ$.
\end{observacion}

\begin{ejemplo}[Distribución estacionaria de una cadena de Markov en CI]\label{ej:estacionaria}
Sea $n=2$ con matriz de transición
\[
K = \begin{pmatrix} 0.7 & 0.3 \\ 0.4 & 0.6 \end{pmatrix}.
\]
La distribución estacionaria satisface $\pi K = \pi$ y $\pi_1+\pi_2=1$:
\[
\pi_1 = 0.7\pi_1 + 0.4\pi_2, \quad \pi_2 = 0.3\pi_1 + 0.6\pi_2.
\]
Resolviendo: $0.3\pi_1 = 0.4\pi_2$, con $\pi_1+\pi_2=1$, se obtiene $\pi = (4/7, 3/7)$. En CI, $\pi$ es el único punto donde $Y^K(\pi)=0$, es decir, el único punto fijo del operador local. La trayectoria desde cualquier $X_0\in\Delta_2^\circ$ satisface $E_\pi(X_t)\searrow 0$ por la Proposición~\ref{prop:disipacion_KL_markov}.
\end{ejemplo}


\section{Sección E — Bloque variacional de alcanzabilidad}

\subsection{Objetivo del bloque}
El objetivo de este bloque es reformular la pregunta
\[
\text{¿Es }P\text{ alcanzable desde }X_0\text{ en }t\text{ pasos bajo acciones admisibles }(\mathcal K_\varepsilon,\mathcal U_R)?
\]
como un problema geométrico y variacional sobre el interior del símplex de probabilidad. La base estructural se toma del marco validado en \cite{ControlInformacional2026}, la geometría afín del símplex y sus arcos elementales se toman de \cite{PistoneShoaib2024Axioms}, y el formalismo del statistical bundle, la acción y las ecuaciones de Euler--Lagrange se toman de \cite{Pistone2018Lagrangian}. El punto de partida informacional es el funcional $E_P$ introducido en la Sección~D (Definición~\ref{def:KL}), cuyo gradiente Fisher y dinámica asociada ya fueron establecidos allí (Proposiciones~\ref{prop:gradKL} y~\ref{prop:solucionExplicita}). Antes de pasar a una formulación variacional, examinaremos si la clase markoviana admisible resulta demasiado amplia y, por tanto, trivializa la pregunta discreta de alcanzabilidad; si eso ocurre, introduciremos una subclase no degenerada para el estudio efectivo. La interpretación local del costo se apoya además en la geometría dual afín y en el teorema de Chentsov, que justifican el uso de la métrica de Fisher como referencia intrínseca de comparación (\cite{Chentsov1982,Amari2016}).

\begin{observacion}[Procedencia de los enunciados]
Las marcas $\fuente{Sección E}$ senalan construcciones propias de esta seccion; las marcas $\fuente{ProInv\_I}$ identifican resultados ya validados en el manuscrito base; y las marcas \texttt{\textbackslash fuentecite\{...\}\{...\}} indican respaldo bibliografico explicito.
\end{observacion}

\subsection{Hipótesis de trabajo}
Fijemos \(n\ge 2\), un objetivo \(P\in\Delta_n^\circ\), un estado inicial \(X_0\in\Delta_n^\circ\) y parámetros \(\varepsilon>0\), \(R>0\), \(\delta\in(0,1)\). En adelante, todo paso discreto se estudiará bajo las restricciones refinadas introducidas más abajo.

\begin{definicion}[Matrices estocásticas por filas \fuente{Sección E}]
Definimos
\[
\mathcal S_n:=\Bigl\{K\in\mathbb R^{n\times n}:\ K_{ij}\ge 0,\ \sum_{j=1}^n K_{ij}=1\text{ para todo }i\Bigr\}.
\]
Los elementos de \(\mathcal S_n\) serán las matrices de transición admisibles para la parte markoviana del problema.
\end{definicion}

\begin{definicion}[Acciones admisibles \fuente{Sección E}]
Definimos
\[
\mathcal K_\varepsilon:=\bigl\{K\in\mathcal S_n:\ K_{ij}\ge \varepsilon\text{ para todo }i,j\bigr\},
\qquad
\mathcal U_R:=\bigl\{v\in\mathbb R^n:\ \|v\|_\infty\le R\bigr\}.
\]
Las acciones de Markov admisibles son los \(K\in\mathcal K_\varepsilon\) y las acciones bayesianas admisibles son las evidencias \(L\) tales que \(\log L\in\mathcal U_R\). Esta clase marca el punto de partida preliminar; la clase de trabajo no degenerada se introducirá más adelante.
\end{definicion}

\begin{lema}[Degeneración de la clase markoviana actual \fuente{Sección E}]
Sea \(Y\in \Delta_n^\circ\) tal que \(Y_i\ge \varepsilon\) para todo \(i\). Definimos la matriz constante por filas
\[
K^Y_{ij}:=Y_j,\qquad 1\le i,j\le n.
\]
Entonces \(K^Y\in \mathcal K_\varepsilon\) y, para todo \(X\in\Delta_n^\circ\), se tiene
\[
L_{K^Y}(X)=XK^Y=Y.
\]
En particular, \(\mathcal K_\varepsilon\) contiene matrices de rango uno.
\end{lema}
\begin{proof}
Para cada \(i\),
\[
\sum_{j=1}^n K^Y_{ij}=\sum_{j=1}^n Y_j=1,
\]
y, por hipotesis, \(K^Y_{ij}=Y_j\ge \varepsilon\). Luego \(K^Y\in\mathcal K_\varepsilon\).

Si \(X\in\Delta_n^\circ\), entonces para cada \(j\),
\[
(L_{K^Y}(X))_j
=
\sum_{i=1}^n X_iK^Y_{ij}
=
\sum_{i=1}^n X_iY_j
=
Y_j\sum_{i=1}^n X_i
=
Y_j.
\]
Por tanto, \(L_{K^Y}(X)=Y\).
\end{proof}

\begin{corolario}[Alcanzabilidad markoviana en un paso \fuente{Sección E}]
Si \(P\in\Delta_{n,\varepsilon}\), entonces \(P\) es alcanzable desde cualquier \(X_0\in\Delta_n^\circ\) en un solo paso mediante una acción de Markov admisible.
\end{corolario}
\begin{proof}
Aplicar el lema anterior con \(Y=P\).
\end{proof}

\begin{observacion}[Necesidad de no degeneración \fuente{Sección E}]
La definición actual de \(\mathcal K_\varepsilon\) no basta, por sí sola, para sostener una teoría discreta no trivial de alcanzabilidad. Para obtener una pregunta genuina de control será necesario imponer alguna restricción adicional, por ejemplo una condición de proximidad a la identidad, una cota espectral uniforme o una hipótesis de irreducibilidad más fuerte.
\end{observacion}

\begin{definicion}[Clase markoviana no degenerada \fuente{Sección E}]
Definimos
\[
\mathcal K_{\varepsilon,\delta}^{\mathrm{nd}}
:=
\left\{
K\in\mathcal S_n:\ 
K_{ij}\ge \varepsilon\ \text{para todo }i,j,
\ \max_{1\le i\le n}\sum_{j=1}^n \bigl|K_{ij}-\delta_{ij}\bigr|\le \delta
\right\},
\]
donde \(\delta_{ij}\) denota la delta de Kronecker. Esta clase excluye el colapso de rango uno y fuerza cada paso markoviano a permanecer a distancia controlada de la identidad.
\end{definicion}

\begin{proposicion}[Paso markoviano acotado \fuente{Sección E}]
Si \(K\in\mathcal K_{\varepsilon,\delta}^{\mathrm{nd}}\), entonces para todo \(X\in\Delta_n\) se tiene
\[
\|L_K(X)-X\|_1\le \delta.
\]
En particular, ninguna acción de Markov en esta clase puede enviar arbitrariamente un estado a otro en un solo paso.
\end{proposicion}

\begin{proof}
Como \(L_K(X)=XK\), se tiene
\[
L_K(X)-X=X(K-I).
\]
Luego,
\[
\|L_K(X)-X\|_1
\le
\sum_{i=1}^n X_i\sum_{j=1}^n |K_{ij}-\delta_{ij}|.
\]
Por la definicion de \(\mathcal K_{\varepsilon,\delta}^{\mathrm{nd}}\),
\[
\sum_{j=1}^n |K_{ij}-\delta_{ij}|\le \delta
\qquad \forall i,
\]
y, como \(\sum_i X_i=1\),
\[
\|L_K(X)-X\|_1\le \delta.
\]
\end{proof}

\begin{definicion}[Núcleo markoviano de suavizado \fuente{Sección E}]
Fijado \(\varepsilon\in(0,1/n)\), definimos la matriz
\[
K_\varepsilon^{\mathrm{sm}}
:=
\bigl(k_{ij}\bigr)_{i,j=1}^n,
\qquad
k_{ij}:=
\begin{cases}
1-(n-1)\varepsilon, & i=j,\\
\varepsilon, & i\neq j.
\end{cases}
\]
Este núcleo es estocástico por filas, tiene todas sus entradas al menos iguales a \(\varepsilon\) y representa una suavización uniforme de proximidad controlada a la identidad.
\end{definicion}

\begin{proposicion}[Propiedades del núcleo de suavizado \fuente{Sección E}]
Si \(0<\varepsilon<1/n\) y \(\delta\ge 2(n-1)\varepsilon\), entonces \(K_\varepsilon^{\mathrm{sm}}\in\mathcal K_{\varepsilon,\delta}^{\mathrm{nd}}\). Además, para todo \(X\in\Delta_n\),
\[
L_{K_\varepsilon^{\mathrm{sm}}}(X)
=(1-n\varepsilon)X+\varepsilon(1,\dots,1).
\]
En particular, \(L_{K_\varepsilon^{\mathrm{sm}}}(X)\in\Delta_n^\circ\) y cada componente queda uniformemente separada de cero por una cota que depende solo de \(\varepsilon\).
\end{proposicion}
\begin{proof}
Para todo \(i\),
\[
\sum_{j=1}^n k_{ij}
=(1-(n-1)\varepsilon)+(n-1)\varepsilon
=1.
\]
Luego \(K_\varepsilon^{\mathrm{sm}}\in\mathcal S_n\) y \(K_{ij}\ge\varepsilon\) para todo \(i,j\). Además, si \(X\in\Delta_n\), entonces para cada \(j\),
\[
(L_{K_\varepsilon^{\mathrm{sm}}}(X))_j
=\sum_{i=1}^n X_i k_{ij}
=(1-(n-1)\varepsilon)X_j+\varepsilon\sum_{i\ne j}X_i
=(1-n\varepsilon)X_j+\varepsilon.
\]
En particular,
\[
\sum_{j=1}^n (L_{K_\varepsilon^{\mathrm{sm}}}(X))_j=1
\quad\text{y}\quad
(L_{K_\varepsilon^{\mathrm{sm}}}(X))_j\ge \varepsilon.
\]
Si \(\delta\ge 2(n-1)\varepsilon\), entonces
\[
\sum_{j=1}^n |k_{ij}-\delta_{ij}|
=2(n-1)\varepsilon\le\delta,
\]
y por tanto \(K_\varepsilon^{\mathrm{sm}}\in\mathcal K_{\varepsilon,\delta}^{\mathrm{nd}}\).
\end{proof}

\begin{definicion}[Familia one-step refinada \fuente{Sección E}]
Asociamos a estas restricciones la familia de pasos compuestos
\[
\mathcal A_{\varepsilon,R,\delta}^{\mathrm{nd}}
:=
\Bigl\{\,T_Y\circ L_K:\ K\in\mathcal K_{\varepsilon,\delta}^{\mathrm{nd}},\ Y\in C_{\mathrm{glob}},\ \|Y\|_\infty\le e^R-1\,\Bigr\},
\]
donde \(L_K(X):=XK\) y \(T_Y\) es el operador informacional global de \cite{ControlInformacional2026}. Equivalentemente, si \(L:=1+Y\), cada paso tiene la forma
\[
X_{t+1}=T_Y(X_tK_t)=B_L(X_tK_t).
\]
Esta es la familia que se usará en adelante para la pregunta de alcanzabilidad no trivial.
\end{definicion}


\begin{definicion}[Trayectoria discreta admisible \fuente{Sección E}]
Sea \(T\in\mathbb N\). Una sucesión \(X_0,\dots,X_T\in\Delta_n^\circ\) se llama \emph{trayectoria discreta admisible} si existe una sucesión \(A_0,\dots,A_{T-1}\) con \(A_t\in\mathcal A_{\varepsilon,R,\delta}^{\mathrm{nd}}\) tal que
\[
X_{t+1}=A_t(X_t),\qquad t=0,\dots,T-1.
\]
\end{definicion}

\begin{definicion}[Símplex truncado \fuente{Sección E}]
Si \(\eta>0\), definimos
\[
\Delta_{n,\eta}:=\bigl\{x\in\Delta_n^\circ:\ x_i\ge \eta\ \text{para todo }i\bigr\}.
\]
\end{definicion}

\begin{lema}[Carta bayesiana local \fuente{Sección E}]
Fijado \(X\in\Delta_n^\circ\), consideramos la aplicación
\[
\Phi_X:\mathcal U_R\to \Delta_n^\circ,
\qquad
\Phi_X(v):=B_{e^v}(X).
\]
Entonces \(\Phi_X\) es de clase \(C^\infty\), satisface \(\Phi_X(0)=X\) y su diferencial en el origen es
\[
D\Phi_X(0)(h)_i
=
X_i\Bigl(h_i-\sum_{j=1}^n X_j h_j\Bigr),
\qquad h\in\mathbb R^n.
\]
En particular, \(\ker D\Phi_X(0)=\operatorname{span}\{(1,\dots,1)\}\) y \(\operatorname{Im}D\Phi_X(0)=T_X\Delta_n^\circ\) (véase Definición~\ref{def:tangente}).
\end{lema}
\begin{proof}
Escribimos
\[
\Phi_X(v)_i=\frac{X_ie^{v_i}}{Z(v)},
\qquad
Z(v):=\sum_{j=1}^n X_je^{v_j}.
\]
Como \(Z(v)>0\), \(\Phi_X\) es de clase \(C^\infty\). Ademas, \(Z(0)=1\), luego \(\Phi_X(0)=X\).

Para \(h\in\mathbb R^n\),
\[
D\Phi_X(0)(h)_i
=
\left.\frac{d}{ds}\right|_{s=0}
\frac{X_ie^{sh_i}}{\sum_{j=1}^n X_je^{sh_j}}
=
X_i\Bigl(h_i-\sum_{j=1}^n X_jh_j\Bigr).
\]
Si \(h=(1,\dots,1)\), entonces \(D\Phi_X(0)(h)=0\). Reciprocamente, si \(D\Phi_X(0)(h)=0\), entonces \(h_i=\sum_j X_jh_j\) para todo \(i\), y por tanto \(h\in\operatorname{span}\{(1,\dots,1)\}\). Luego
\[
\ker D\Phi_X(0)=\operatorname{span}\{(1,\dots,1)\}.
\]
Ademas,
\[
\sum_{i=1}^n D\Phi_X(0)(h)_i
=
\sum_{i=1}^n X_ih_i-\sum_{j=1}^n X_jh_j\sum_{i=1}^n X_i
=
0,
\]
por lo que \(\operatorname{Im}D\Phi_X(0)\subset T_X\Delta_n^\circ\). Como ambos espacios tienen dimension \(n-1\), la inclusion es igualdad.
\end{proof}

\begin{corolario}[Inversion local de la carta bayesiana \fuente{Sección E}]
Sea
\[
H_X:=\bigl\{h\in\mathbb R^n:\ \sum_{j=1}^n X_j h_j=0\bigr\}.
\]
La restriccion \(\Phi_X|_{H_X}\) es un difeomorfismo local de clase \(C^\infty\) entre un entorno de \(0\in H_X\) y un entorno de \(X\) en \(\Delta_n^\circ\). En particular, todo estado suficientemente cercano a \(X\) se escribe de manera unica como \(B_{e^v}(X)\) con \(v\in H_X\) pequeno.
\end{corolario}
\begin{proof}
Sea
\[
H_X=\bigl\{h\in\mathbb R^n:\ \sum_{j=1}^n X_jh_j=0\bigr\}.
\]
Si \(h\in H_X\), entonces
\[
D\Phi_X(0)(h)_i=X_ih_i.
\]
Por tanto, \(D\Phi_X(0)|_{H_X}\) es inyectiva. Como
\[
\dim H_X=n-1=\dim T_X\Delta_n^\circ,
\]
es un isomorfismo lineal. El teorema de la funcion inversa aplica a \(\Phi_X|_{H_X}\), y la conclusion sigue.
\end{proof}

\begin{proposicion}[Control multiplicativo bajo evidencia acotada \fuente{Sección E}]
Sea \(L\in\mathbb R_{>0}^n\) tal que \(\log L\in\mathcal U_R\). Entonces, para todo \(X\in\Delta_n^\circ\) y todo \(i=1,\dots,n\),
\[
e^{-2R}X_i\le B_L(X)_i\le e^{2R}X_i.
\]
En particular,
\[
B_L(X)\in \Delta_{n,\eta_R(X)},
\qquad
\eta_R(X):=e^{-2R}\min_{1\le i\le n}X_i.
\]
\end{proposicion}
\begin{proof}
De \(\log L\in\mathcal U_R\) se sigue
\[
e^{-R}\le L_i\le e^R,\qquad i=1,\dots,n.
\]
Si \(c:=\sum_{j=1}^n X_jL_j\), entonces
\[
e^{-R}\le c\le e^R.
\]
Por definicion,
\[
B_L(X)_i=\frac{X_iL_i}{c}.
\]
Luego,
\[
\frac{e^{-R}}{e^R}X_i\le B_L(X)_i\le \frac{e^R}{e^{-R}}X_i,
\]
esto es,
\[
e^{-2R}X_i\le B_L(X)_i\le e^{2R}X_i.
\]
La cota inferior implica
\[
B_L(X)_i\ge e^{-2R}\min_{1\le j\le n}X_j=\eta_R(X).
\]
\end{proof}

\begin{proposicion}[Identificacion con la correccion multiplicativa \fuentecite{ProInv\_I}{ControlInformacional2026}]
Sea \(L\in\mathbb R_{>0}^n\) y defínase \(Y:=L-1\). Entonces, para todo \(X\in\Delta_n^\circ\),
\[
B_L(X)=T_Y(X),
\]
es decir, la actualización bayesiana usada aquí coincide con la corrección multiplicativa global introducida en \cite{ControlInformacional2026} bajo el cambio de variable \(L=1+Y\).
\end{proposicion}
\begin{proof}
La fórmula de \(T_Y\) establecida en la Definición~\ref{def:Tglobal} de la Sección~C es
\[
T_Y(X)_i=\frac{X_i(1+Y_i)}{\sum_{j=1}^n X_j(1+Y_j)}.
\]
Bajo el cambio de variable \(L := 1+Y\), esta expresión coincide exactamente con \(B_L(X)_i = \frac{X_iL_i}{\sum_{j=1}^n X_jL_j}\).
\end{proof}

\begin{proposicion}[Cota del campo de correccion \fuente{Sección E}]
Si \(\log L\in\mathcal U_R\) y \(Y:=L-1\), entonces \(\|Y\|_\infty\le e^R-1\). En particular, la correccion multiplicativa global \(T_Y\) usada aqui actua con amplitud uniformemente acotada por \(R\).
\end{proposicion}
\begin{proof}
De \(\log L\in\mathcal U_R\) se sigue
\[
e^{-R}\le L_i\le e^R,
\qquad i=1,\dots,n.
\]
Si \(Y:=L-1\), entonces
\[
|Y_i|=|L_i-1|\le e^R-1.
\]
Luego
\[
\|Y\|_\infty\le e^R-1.
\]
\end{proof}

\begin{proposicion}[Cota del paso compuesto \fuente{Sección E}]
Sea \(X\in\Delta_n^\circ\), \(K\in\mathcal K_{\varepsilon,\delta}^{\mathrm{nd}}\) y \(L\in\mathbb R_{>0}^n\) con \(\log L\in\mathcal U_R\). Entonces
\[
\|B_L(XK)-X\|_1\le \delta + e^{2R}-1.
\]
Si ademas \(X\in\Delta_{n,\eta}\) para algún \(\eta>0\), entonces
\[
g_X\bigl(B_L(XK)-X,\ B_L(XK)-X\bigr)\le \eta^{-1}\bigl(\delta + e^{2R}-1\bigr)^2.
\]
\end{proposicion}
\begin{proof}
Sea \(Y:=XK\). Por la proposición de paso markoviano acotado,
\[
\|Y-X\|_1\le \delta.
\]
Además, si \(c:=\sum_{j=1}^n Y_jL_j\), entonces
\[
e^{-R}\le c\le e^R,
\qquad
e^{-R}\le L_i\le e^R.
\]
Por tanto,
\[
e^{-2R}\le \frac{L_i}{c}\le e^{2R},
\qquad i=1,\dots,n.
\]
Luego
\[
|B_L(Y)_i-Y_i|
=Y_i\left|\frac{L_i}{c}-1\right|
\le Y_i\,(e^{2R}-1),
\]
y al sumar en \(i\) se obtiene
\[
\|B_L(Y)-Y\|_1\le e^{2R}-1.
\]
Finalmente,
\[
\|B_L(XK)-X\|_1
\le \|B_L(Y)-Y\|_1+\|Y-X\|_1
\le \delta+e^{2R}-1.
\]
Si además \(X\in\Delta_{n,\eta}\), entonces para \(u:=B_L(XK)-X\),
\[
g_X(u,u)=\sum_{i=1}^n \frac{u_i^2}{X_i}
\le \eta^{-1}\sum_{i=1}^n u_i^2
\le \eta^{-1}\left(\sum_{i=1}^n |u_i|\right)^2
=\eta^{-1}\|u\|_1^2.
\]
\end{proof}

\begin{observacion}[Comparabilidad de la métrica de Fisher]
La definición del costo elemental fija el punto base en \(X\); por ello la evaluación en \(g_X\) es la adecuada. En subconjuntos compactos \(\Delta_{n,\eta}\), las métricas de Fisher son uniformemente equivalentes, de modo que medir en el extremo inicial o en el extremo final produce cotas comparables hasta una constante uniforme.
\end{observacion}
\begin{definicion}[Costo elemental discreto \fuente{Sección E}]
Sea \(X\in\Delta_n^\circ\), \(K\in\mathcal K_{\varepsilon,\delta}^{\mathrm{nd}}\) y \(L\in\mathbb R_{>0}^n\) con \(\log L\in\mathcal U_R\). Definimos el costo elemental del paso compuesto \(X\mapsto B_L(XK)\) por
\[
\mathfrak c_{P,\lambda}(X,K,L):=\frac12\,g_X\bigl(B_L(XK)-X,\ B_L(XK)-X\bigr)+\lambda\,E_P\bigl(B_L(XK)\bigr).
\]
Aqui \(g_X\) es la métrica de Fisher de la Definición~\ref{def:Fisher} y \(E_P\) es el funcional de la Definición~\ref{def:KL}.
Si \((X_t,K_t,L_t)_{t=0}^{T-1}\) es una trayectoria discreta admisible, su costo total es
\[
\mathfrak A_{P,\lambda}:=\sum_{t=0}^{T-1}\mathfrak c_{P,\lambda}(X_t,K_t,L_t).
\]
\end{definicion}

\begin{observacion}[Lagrangiano como relajacion continua \fuente{Sección E}]
El funcional \(\mathcal A_P\) se interpreta aqui como la relajacion formal de \(\mathfrak A_{P,\lambda}\) cuando el mallado temporal se refina y los incrementos discretos son de primer orden. En este bloque no se afirma una teoria completa de \(\Gamma\)-convergencia; solo se fija el funcional continuo que servira como blanco variacional.
\end{observacion}

\begin{proposicion}[Persistencia en el interior \fuente{Sección E}]
Sea \(T\in\mathbb N\) y sea \(X_0,\dots,X_T\) una trayectoria discreta admisible. Entonces existe
\[
 \eta_*:=e^{-2R}\min\!\bigl\{\varepsilon,\min_{1\le i\le n}(X_0)_i\bigr\}>0
\]
tal que \(X_t\in \Delta_{n,\eta_*}\) para todo \(t=0,\dots,T\). En particular, la imagen de la trayectoria queda contenida en un subconjunto compacto de \(\Delta_n^\circ\), con cota uniforme independiente de \(T\).
\end{proposicion}
\begin{proof}
Escribamos
\[
X_{t+1}=B_{L_t}(X_tK_t).
\]
Para cada \(i\),
\[
(X_tK_t)_i=\sum_{j=1}^n (X_t)_j (K_t)_{ji}
\ge \varepsilon\sum_{j=1}^n (X_t)_j
=\varepsilon.
\]
Por la proposición de control multiplicativo,
\[
(X_{t+1})_i\ge e^{-2R}(X_tK_t)_i\ge e^{-2R}\varepsilon.
\]
Como \(X_0\in\Delta_n^\circ\), se tiene también \((X_0)_i\ge \min_j(X_0)_j\). En consecuencia, si
\[
\eta_*:=e^{-2R}\min\Bigl\{\varepsilon,\min_{1\le j\le n}(X_0)_j\Bigr\},
\]
entonces \(X_t\in\Delta_{n,\eta_*}\) para todo \(t=0,\dots,T\).
\end{proof}

\begin{definicion}[Arcos geométricos elementales \fuentecite{ref1}{PistoneShoaib2024Axioms}]
Para \(p,q\in\Delta_n^\circ\), definimos el \emph{arco de mezcla} y el \emph{arco exponencial} por
\[
m_{p,q}(s):=(1-s)p+sq,
\qquad
e_{p,q}(s)_i:=\frac{p_i^{1-s}q_i^s}{\sum_{j=1}^n p_j^{1-s}q_j^s},
\qquad s\in[0,1].
\]
\end{definicion}

\begin{proposicion}[Regularidad de los arcos \fuentecite{ref1}{PistoneShoaib2024Axioms}]
Para cualesquiera \(p,q\in\Delta_n^\circ\), las aplicaciones \(m_{p,q}\) y \(e_{p,q}\) son curvas \(C^\infty\) en \([0,1]\) con valores en \(\Delta_n^\circ\), y satisfacen
\[
m_{p,q}(0)=e_{p,q}(0)=p,
\qquad
m_{p,q}(1)=e_{p,q}(1)=q.
\]
\end{proposicion}
\begin{proof}
El arco de mezcla
\[
m_{p,q}(s)=(1-s)p+sq
\]
es afín en \(s\); por tanto, es \(C^\infty\). Para el arco exponencial, si
\[
D(s):=\sum_{j=1}^n p_j^{1-s}q_j^s,
\]
entonces \(D(s)>0\) y \(D\in C^\infty([0,1])\). Cada componente
\[
e_{p,q}(s)_i=\frac{p_i^{1-s}q_i^s}{D(s)}
\]
es cociente de funciones \(C^\infty\) con denominador positivo; luego \(e_{p,q}\in C^\infty([0,1];\Delta_n^\circ)\). Además,
\[
m_{p,q}(0)=e_{p,q}(0)=p,
\qquad
m_{p,q}(1)=e_{p,q}(1)=q.
\]
\end{proof}

\begin{observacion}[Interpolación continua auxiliar \fuente{Sección E}]
Toda trayectoria discreta admisible \(X_0,\dots,X_T\) admite una interpolación continua por concatenación de arcos de mezcla o de arcos exponenciales. Esta interpolación se usa solo como representante geométrico de la sucesión discreta; no se afirma ninguna propiedad minimizante para la acción.
\end{observacion}

\begin{definicion}[Función potencial de referencia \fuentecite{ProInv\_I}{ControlInformacional2026}]
Para \(P\in\Delta_n^\circ\), definimos
\[
E_P(X):=D_{\mathrm{KL}}(P\|X)
=\sum_{i=1}^n P_i\log\!\left(\frac{P_i}{X_i}\right),
\qquad X\in\Delta_n^\circ.
\]
\end{definicion}

\begin{proposicion}[Propiedades del potencial KL \fuentecite{ProInv\_I}{ControlInformacional2026}]
La función \(E_P:\Delta_n^\circ\to\mathbb R\) es de clase \(C^\infty\), convexa en toda carta afín del símplex, y posee un único minimizador global, dado por \(X=P\). Además, su gradiente respecto de la métrica de Fisher satisface
\[
\operatorname{grad}_g E_P(X)=X-P,
\qquad X\in\Delta_n^\circ.
\]
\end{proposicion}
\begin{proof}
Se tiene
\[
E_P(X)=\sum_{i=1}^n P_i\log P_i-\sum_{i=1}^n P_i\log X_i.
\]
Por tanto \(E_P\in C^\infty(\Delta_n^\circ)\). Para \(u\in T_X\Delta_n^\circ\), \(\sum_i u_i=0\), y
\[
DE_P(X)[u]=-\sum_{i=1}^n \frac{P_i}{X_i}u_i.
\]
Si \(g_X(u,v)=\sum_i u_iv_i/X_i\), entonces
\[
g_X(X-P,u)=\sum_{i=1}^n \frac{X_i-P_i}{X_i}u_i
=\sum_{i=1}^n u_i-\sum_{i=1}^n \frac{P_i}{X_i}u_i
=-\sum_{i=1}^n \frac{P_i}{X_i}u_i.
\]
Luego
\[
\operatorname{grad}_g E_P(X)=X-P.
\]
Además,
\[
D^2E_P(X)[u,u]=\sum_{i=1}^n \frac{P_i}{X_i^2}u_i^2>0
\qquad(u\neq 0,\ \sum_i u_i=0),
\]
por lo que \(E_P\) es estrictamente convexa en cada carta afín. Finalmente,
\[
E_P(X)=D_{\mathrm{KL}}(P\|X)\ge 0,
\qquad
E_P(X)=0\iff X=P,
\]
y por tanto \(P\) es el único minimizador global.
\end{proof}

\begin{proposicion}[Cuadraticidad local del potencial \fuentecite{Sección E}{ControlInformacional2026}]
El diferencial de \(E_P\) se anula en \(P\) y el hessiano de \(E_P\) en \(P\) coincide con la métrica de Fisher. En particular, en coordenadas tangentes centradas en \(P\),
\[
E_P(P+\xi)=\frac12\,g_P(\xi,\xi)+o(\|\xi\|^2)
\qquad (\xi\to 0,\ \sum_i \xi_i=0).
\]
\end{proposicion}
\begin{proof}
Por la proposición anterior, \(DE_P(P)=0\). Además, para \(\xi\in T_P\Delta_n^\circ\),
\[
D^2E_P(P)[\xi,\xi]
=\sum_{i=1}^n \frac{P_i}{P_i^2}\,\xi_i^2
=\sum_{i=1}^n \frac{\xi_i^2}{P_i}
=g_P(\xi,\xi).
\]
El desarrollo de Taylor en una carta afín centrada en \(P\) da
\[
E_P(P+\xi)
=E_P(P)+DE_P(P)[\xi]+\frac12 D^2E_P(P)[\xi,\xi]+o(\|\xi\|^2)
=\frac12 g_P(\xi,\xi)+o(\|\xi\|^2).
\]
\end{proof}

\begin{definicion}[Statistical bundle finito \fuentecite{ref2}{Pistone2018Lagrangian}]
Denotamos por
\[
S\Delta_n^\circ:=\{(q,w)\in \Delta_n^\circ\times\mathbb R^n:\ \sum_{i=1}^n w_i=0\}
\]
el statistical bundle finito. En cada fibra \(T_q\Delta_n^\circ\) se considera la forma bilineal de Fisher
\[
g_q(u,v):=\sum_{i=1}^n \frac{u_i v_i}{q_i}.
\]
\end{definicion}

\begin{definicion}[Lagrangiano de alcanzabilidad \fuentecite{ref2}{Pistone2018Lagrangian}]
Fijado \(\lambda>0\), definimos el Lagrangiano
\[
\mathcal L_P(q,w):=\frac12\,g_q(w,w)+\lambda\,E_P(q),
\qquad (q,w)\in S\Delta_n^\circ.
\]
Si \(q:[0,T]\to \Delta_n^\circ\) es una curva de clase \(C^1\), su acción es
\[
\mathcal A_P(q):=\int_0^T \mathcal L_P\bigl(q(t),\dot q(t)\bigr)\,dt.
\]
\end{definicion}

\begin{proposicion}[Levantamiento al statistical bundle \fuentecite{ref2}{Pistone2018Lagrangian}]
Si \(q:[0,T]\to\Delta_n^\circ\) es de clase \(C^2\), entonces
\[
\gamma(t):=\bigl(q(t),\dot q(t)\bigr)
\]
define una curva de clase \(C^1\) en \(S\Delta_n^\circ\). En particular, la velocidad y la aceleración de \(q\) quedan bien definidas en el marco afín del statistical bundle descrito en \cite{Pistone2018Lagrangian}.
\end{proposicion}
\begin{proof}
Si \(q\in C^2([0,T];\Delta_n^\circ)\), entonces \(\dot q\in C^1([0,T];\mathbb R^n)\). Como
\[
\sum_{i=1}^n q_i(t)=1,
\]
al derivar se obtiene
\[
\sum_{i=1}^n \dot q_i(t)=0.
\]
Por tanto,
\[
\gamma(t):=(q(t),\dot q(t))\in S\Delta_n^\circ
\qquad\text{para todo }t\in[0,T],
\]
y \(\gamma\in C^1([0,T];S\Delta_n^\circ)\).
\end{proof}

\begin{teorema}[Existencia de minimizador para el problema relajado \fuente{Sección E}]
Sean \(X_0,P\in\Delta_n^\circ\) y \(T>0\). Consideremos la clase de curvas
\[
\mathcal C_{\eta_T}(X_0,P;T)
:=
\bigl\{q\in H^1([0,T];\Delta_n^\circ):\ q(0)=X_0,\ q(T)=P,\ q([0,T])\subset \Delta_{n,\eta_T}\bigr\}.
\]
Entonces la acción \(\mathcal A_P\) es acotada inferiormente en \(\mathcal C_{\eta_T}(X_0,P;T)\) y, bajo la coercividad del término cinético, admite al menos un minimizador en dicha clase.
\end{teorema}
\begin{proof}
Sea \((q^m)_{m\in\mathbb N}\subset \mathcal C_{\eta_T}(X_0,P;T)\) una sucesión minimizante. Como \(q^m([0,T])\subset\Delta_{n,\eta_T}\) y \(\Delta_{n,\eta_T}\) es compacto, el término potencial \(E_P(q^m)\) es uniformemente acotado inferiormente. La coercividad del término cinético implica acotación de \((\dot q^m)\) en \(L^2([0,T])\); por tanto, \((q^m)\) es acotada en \(H^1([0,T])\).

Por reflexividad, existe una subsucesión, todavía denotada \(q^m\), y una curva \(q^*\in H^1([0,T])\) tales que
\[
q^m\rightharpoonup q^* \text{ en } H^1,
\qquad
q^m\to q^* \text{ en } C^0([0,T]).
\]
Luego \(q^*(0)=X_0\), \(q^*(T)=P\) y \(q^*([0,T])\subset\Delta_{n,\eta_T}\). Por semicontinuidad inferior del término cinético y continuidad de \(E_P\) respecto de la convergencia uniforme,
\[
\mathcal A_P(q^*)\le \liminf_{m\to\infty}\mathcal A_P(q^m)=\inf_{q\in\mathcal C_{\eta_T}(X_0,P;T)}\mathcal A_P(q).
\]
Por consiguiente, \(q^*\) es minimizador.
\end{proof}

\begin{proposicion}[Ecuaciones de Euler--Lagrange \fuentecite{ref2}{Pistone2018Lagrangian}]
Si \(q\in \mathcal C_{\eta_T}(X_0,P;T)\) es un minimizador de \(\mathcal A_P\) y \(q\) es de clase \(C^2\), entonces satisface la ecuación de Euler--Lagrange
\[
\frac{d}{dt}\bigl(\partial_w \mathcal L_P(q(t),\dot q(t))\bigr)
-\partial_q \mathcal L_P(q(t),\dot q(t))=0
\]
en cualquier carta del statistical bundle.
\end{proposicion}
\begin{proof}
Sea \(\varphi\in C_c^\infty((0,T);\mathbb R^n)\) tal que \(\sum_{i=1}^n \varphi_i=0\). Para \(q^\epsilon:=q+\epsilon\varphi\),
\[
0=\frac{d}{d\epsilon}\Big|_{\epsilon=0}\mathcal A_P(q^\epsilon)
=\int_0^T \sum_{i=1}^n\left(
\frac{\dot q_i}{q_i}\dot\varphi_i
-\frac12\frac{\dot q_i^2}{q_i^2}\varphi_i
-\lambda\frac{P_i}{q_i}\varphi_i
\right)dt.
\]
Integrando por partes y usando \(\varphi(0)=\varphi(T)=0\), queda
\[
\int_0^T \sum_{i=1}^n \left(
-\frac{d}{dt}\Bigl(\frac{\dot q_i}{q_i}\Bigr)
-\frac12\frac{\dot q_i^2}{q_i^2}
-\lambda\frac{P_i}{q_i}
\right)\varphi_i\,dt=0.
\]
El anulador del subespacio \(\{\varphi:\sum_i\varphi_i=0\}\) es \(\operatorname{span}\{(1,\dots,1)\}\); por tanto existe \(\mu(t)\) tal que
\[
\frac{d}{dt}\Bigl(\partial_w\mathcal L_P(q(t),\dot q(t))\Bigr)-\partial_q\mathcal L_P(q(t),\dot q(t))=0
\]
en cualquier carta del statistical bundle.
\end{proof}

\begin{corolario}[Caso puramente cinético \textnormal{(ref1+ref2; \citealp{PistoneShoaib2024Axioms,Pistone2018Lagrangian})}]
Si \(\lambda=0\), entonces las trayectorias críticas de \(\mathcal A_P\) son geodésicas de la geometría afín del símplex. En las cartas exponenciales y de mezcla de \cite{PistoneShoaib2024Axioms}, estas trayectorias son precisamente las curvas de aceleración nula.
\end{corolario}
\begin{proof}
Si \(\lambda=0\), entonces
\[
\mathcal L_P(q,w)=\frac12\,g_q(w,w).
\]
Las curvas críticas de la energía cinética en una variedad riemanniana satisfacen la ecuación geodésica. En las cartas de mezcla y exponenciales de \cite{PistoneShoaib2024Axioms}, dichas geodésicas son exactamente las curvas de aceleración nula. Luego las trayectorias críticas de \(\mathcal A_P\) son geodésicas de la geometría afín del símplex.
\end{proof}

\begin{teorema}[Discretización admisible del minimizador \fuente{Sección E}]
Sea \(q^*\in \mathcal C_{\eta_T}(X_0,P;T)\) un minimizador de \(\mathcal A_P\) de clase \(C^2\). Fijemos el núcleo markoviano de suavizado \(K_\varepsilon^{\mathrm{sm}}\) de la Definición anterior, con \(0<\varepsilon<1/n\) y \(\delta\ge 2(n-1)\varepsilon\).

Entonces existe una partición \(0=t_0<t_1<\cdots<t_m=T\) suficientemente fina y existen constantes \(h_0>0\) y \(R_0>0\) tales que, si \(h:=\max_k(t_{k+1}-t_k)\le h_0\) y definimos
\[
z_k:=q^*(t_k)K_\varepsilon^{\mathrm{sm}},\qquad
L_{k,i}:=\frac{q_i^*(t_{k+1})}{(z_k)_i},\qquad i=1,\dots,n,
\]
entonces \(L_k=(L_{k,1},\dots,L_{k,n})\in\mathbb R_{>0}^n\), \(\log L_k\in\mathcal U_{R_0}\) y
\[
q^*(t_{k+1})=B_{L_k}(z_k)=B_{L_k}\bigl(q^*(t_k)K_\varepsilon^{\mathrm{sm}}\bigr),\qquad k=0,\dots,m-1.
\]
En particular, para todo \(R\ge R_0\) la sucesión \(q^*(t_0),\dots,q^*(t_m)\) es una trayectoria discreta admisible bajo \(\mathcal A_{\varepsilon,R,\delta}^{\mathrm{nd}}\), realizada mediante el paso compuesto canónico
\[
X_{t+1}=B_{L_t}\bigl(X_tK_\varepsilon^{\mathrm{sm}}\bigr).
\]
\end{teorema}
\begin{proof}
Sea \(q^*\in\mathcal C_{\eta_T}(X_0,P;T)\). Fijemos una partición \(0=t_0<t_1<\cdots<t_m=T\) y pongamos
\[
z_k:=q^*(t_k)K_\varepsilon^{\mathrm{sm}},
\qquad
L_{k,i}:=\frac{q_i^*(t_{k+1})}{(z_k)_i}.
\]
Como \(q^*(t_k)\in\Delta_{n,\eta_T}\) y cada entrada de \(K_\varepsilon^{\mathrm{sm}}\) es al menos \(\varepsilon\),
\[
(z_k)_i=\sum_{j=1}^n q_j^*(t_k)(K_\varepsilon^{\mathrm{sm}})_{ji}\ge \varepsilon.
\]
Además, \(q_i^*(t_{k+1})\in[\eta_T,1]\), luego
\[
\eta_T\le L_{k,i}\le \frac1\varepsilon.
\]
Por consiguiente,
\[
\log L_k\in\mathcal U_{R_0},
\qquad
R_0:=\max\{-\log\eta_T,\log(1/\varepsilon)\}.
\]
Ahora,
\[
\sum_{j=1}^n (z_k)_jL_{k,j}
=\sum_{j=1}^n q_j^*(t_{k+1})=1,
\]
y por tanto
\[
B_{L_k}(z_k)_i
=\frac{(z_k)_iL_{k,i}}{\sum_{j=1}^n (z_k)_jL_{k,j}}
=q_i^*(t_{k+1}).
\]
De aquí
\[
q^*(t_{k+1})=B_{L_k}\bigl(q^*(t_k)K_\varepsilon^{\mathrm{sm}}\bigr)
\]
para todo \(k\). Si \(R\ge R_0\), entonces cada paso pertenece a \(\mathcal A_{\varepsilon,R,\delta}^{\mathrm{nd}}\).
\end{proof}

\begin{proposicion}[Euler--Lagrange en coordenadas baricentricas \fuente{Sección E}]
Sea \(q\) una curva critica de \(\mathcal A_P\) con \(\sum_{i=1}^n q_i(t)=1\). Entonces existe una funcion escalar \(\mu(t)\) tal que, para cada \(i=1,\dots,n\),
\[
\frac{d}{dt}\left(\frac{\dot q_i(t)}{q_i(t)}\right)+\frac12\,\frac{\dot q_i(t)^2}{q_i(t)^2}+\lambda\,\frac{P_i}{q_i(t)}=\mu(t),
\]
equivalentemente,
\[
\ddot q_i(t)-\frac12\,\frac{\dot q_i(t)^2}{q_i(t)}+\lambda P_i-\mu(t)\,q_i(t)=0.
\]
La funcion \(\mu(t)\) es el multiplicador de Lagrange asociado a la restriccion \(\sum_i q_i=1\).
\end{proposicion}
\begin{proof}
Consideremos el funcional aumentado
\[
\widetilde{\mathcal L}_P(q,\dot q,\mu)
:=\frac12\sum_{i=1}^n\frac{\dot q_i^2}{q_i}
+\lambda\sum_{i=1}^n P_i\log\!\left(\frac{P_i}{q_i}\right)
-\mu(t)\Bigl(\sum_{i=1}^n q_i-1\Bigr).
\]
Entonces
\[
\partial_{\dot q_i}\widetilde{\mathcal L}_P=\frac{\dot q_i}{q_i},
\qquad
\partial_{q_i}\widetilde{\mathcal L}_P
=-\frac12\frac{\dot q_i^2}{q_i^2}-\lambda\frac{P_i}{q_i}-\mu(t).
\]
La ecuación de Euler--Lagrange da
\[
\frac{d}{dt}\left(\frac{\dot q_i}{q_i}\right)
+\frac12\frac{\dot q_i^2}{q_i^2}
+\lambda\frac{P_i}{q_i}
=\mu(t),
\qquad i=1,\dots,n.
\]
Multiplicando por \(q_i\) y usando
\[
q_i\frac{d}{dt}\left(\frac{\dot q_i}{q_i}\right)=\ddot q_i-\frac{\dot q_i^2}{q_i},
\]
se obtiene
\[
\ddot q_i-\frac12\frac{\dot q_i^2}{q_i}+\lambda P_i-\mu(t)q_i=0.
\]
\end{proof}

\begin{corolario}[Límite estacionario formal \fuente{Sección E}]
Consideremos el funcional reescalado \(\mathcal L_{P,\eta}(q,w):=\frac{\eta}{2}\,g_q(w,w)+\lambda E_P(q)\). Al pasar formalmente al límite \(\eta\downarrow 0\), las trayectorias críticas quedan forzadas a satisfacer
\[
\operatorname{grad}_g E_P(q)=0,
\]
y por tanto
\[
q=P.
\]
Este paso no produce por sí mismo un flujo de primer orden; sólo produce la condición estacionaria del potencial.
\end{corolario}
\begin{proof}
La ecuación variacional asociada a \(\mathcal L_{P,\eta}\) contiene el término inercial multiplicado por \(\eta\) y el término de potencial \(\lambda\,\operatorname{grad}_gE_P(q)\). Al pasar formalmente al límite \(\eta\downarrow 0\), desaparece la contribución inercial y queda
\[
\operatorname{grad}_gE_P(q)=0.
\]
Por la proposición sobre el potencial KL,
\[
\operatorname{grad}_gE_P(q)=q-P.
\]
Luego \(q=P\).
\end{proof}

\begin{observacion}[Flujo heredado de la Sección~D]
El flujo \(\dot X = P - X\) y su convergencia a \(P\) son consecuencia directa de la Proposición~\ref{prop:gradKL} y la Proposición~\ref{prop:solucionExplicita} de la Sección~D. Este resultado no se rededuce aquí; se adopta como punto de partida del análisis variacional.
\end{observacion}

\subsection{Observaciones geométricas finales}

\begin{observacion}[Dualidad \(e/m\) del paso compuesto \fuentecite{ref3}{Amari2016}]
El paso libre \(X_tK_\varepsilon^{\mathrm{sm}}\) y la corrección bayesiana \(B_{L_t}\) se interpretan, respectivamente, como una componente \(m\)-afín y una componente \(e\)-afín del proceso compuesto. Bajo esta lectura, la función de costo separa el transporte libre y la corrección informacional, y el caso \(\lambda=0\) del Corolario anterior corresponde al régimen puramente geodésico del símplex dually flat.
\end{observacion}

\begin{observacion}[Fisher como métrica natural \fuentecite{ref3}{Amari2016,Chentsov1982}]
La selección de la métrica de Fisher en el funcional continuo no es arbitraria: por el teorema de Chentsov, esta es la métrica intrínseca compatible con las transformaciones markovianas finitas. En consecuencia, el puente discreto-continuo debe formularse en términos de Fisher/KL y no mediante una métrica composicional auxiliar.
\end{observacion}

\begin{problema}[Criterio variacional de alcanzabilidad \fuente{Sección E}]
Determinar condiciones necesarias y suficientes para que, dados \(X_0,P\in \Delta_n^\circ\) y \(T\in\mathbb N\), exista una trayectoria discreta admisible de longitud \(T\) que conecte \(X_0\) con \(P\) bajo las acciones refinadas \((\mathcal K_{\varepsilon,\delta}^{\mathrm{nd}},\mathcal U_R)\), y comparar dicho criterio con la existencia de minimizadores de la acción relajada \(\mathcal A_P\).
\end{problema}


\section{Conclusiones y perspectivas}

En este trabajo se desarrolló un marco interno para el estudio de transformaciones entre distribuciones de probabilidad discretas sobre el símplex. En particular, se introdujo el concepto de correcciones multiplicativas y su estructura asociada, distinguiendo entre correcciones locales dependientes del estado y correcciones globales definidas de manera uniforme. A partir de ello, se construyeron operadores informacionales que preservan el símplex y cuya ley de composición corresponde a un producto coordenado de factores positivos.

Asimismo, se estableció una conexión con la geometría de información mediante la métrica de Fisher y la divergencia de Kullback--Leibler, mostrando que la dinámica inducida por el gradiente conduce a un flujo explícito que converge hacia una distribución de referencia. Esto proporciona una interpretación dinámica del proceso de corrección como una relajación informacional.

En continuidad con este desarrollo, el manuscrito unificadoI podría orientarse hacia diversas extensiones estructurales del marco descrito. En particular, resulta natural explorar la relación entre los operadores introducidos y operadores de tipo Markov, con el objetivo de analizar si estos últimos pueden interpretarse como una subfamilia dentro del conjunto de transformaciones multiplicativas. De manera análoga, podría estudiarse la existencia de una representación multiplicativa de operadores de Markov y la posible emergencia de una estructura de tipo grupoide.

Por otra parte, una dirección relevante consiste en incorporar el enfoque del análisis de datos composicionales de Aitchison, en el cual el símplex se dota de una geometría basada en coordenadas log-ratio. En este contexto, podría investigarse si las correcciones multiplicativas del marco descrito admiten una representación lineal en dichas coordenadas, lo que permitiría reinterpretar las transformaciones en términos de estructuras vectoriales asociadas y analizar nociones como ortogonalidad y proyección.

Finalmente, también podría considerarse la conexión con problemas de inferencia estadística. En particular, sería de interés analizar si ciertos procedimientos, como la actualización bayesiana, pueden interpretarse dentro de este esquema como correcciones óptimas, así como explorar posibles vínculos con herramientas del análisis multivariado composicional, incluyendo matrices de covarianza en coordenadas log-ratio, métodos tipo PCA y técnicas de agrupamiento. Estas líneas constituyen posibles direcciones de desarrollo y requerirán un análisis más detallado en etapas posteriores.
\clearpage


\bibliography{Manuscrito_Unificado}
\end{document}

